%% GENERATED FILE -- do not hand-edit.
%% SOURCE: experiments/database/scripts/build_hosts_candidate_datasets_table.py
\begin{landscape}
\begingroup
\footnotesize
\setlength{\tabcolsep}{4pt}
\renewcommand{\arraystretch}{1.15}
\begin{longtable}{@{}r@{\hspace{4pt}} L{78mm} L{72mm} L{76mm}@{}}
\caption[]{Sources for the whole pool of Table~\ref{tab:candidates}, in the same order.
\emph{Why it is ordered here} is the band reason for a host-bootstrap, regulatory-DNA or
cross-host-fitness row, and the tier and measurement rule for a scale row. \emph{Where
the per-record data lives} is the deposit, with its retrieval class: a plain entry is
fetched programmatically, \emph{browser} needs a human because the host returns 403 or
renders in JavaScript, and \emph{ask} means no public copy exists and the request is in
Table~\ref{tab:asks}. Every link is clickable.}
\label{tab:sources-all}\\
\toprule
\textbf{\#} & \textbf{Dataset, citation and link} & \textbf{Why it is ordered here} & \textbf{Where the per-record data lives} \\
\midrule
\endfirsthead
\multicolumn{4}{@{}l}{\footnotesize\emph{Table~\ref{tab:sources-all}, continued}}\\
\toprule
\textbf{\#} & \textbf{Dataset, citation and link} & \textbf{Why it is ordered here} & \textbf{Where the per-record data lives} \\
\midrule
\endhead
\bottomrule
\endfoot
1 & \textbf{de Boer 2020 (GPRA sort-seq, 100M random promoters)} \org{S. cerevisiae}\newline de Boer CG, Vaishnav ED, Sadeh R, Abeyta EL, Friedman N, Regev A. Nat Biotechnol 2020;38:56-65.\newline \href{https://doi.org/10.1038/s41587-019-0315-8}{doi:\allowbreak 10.\allowbreak 1038/\allowbreak s41587-\allowbreak 019-\allowbreak 0315-\allowbreak 8} & Regulatory DNA band: one row carrying more measurements than every other row in this wave combined, on a schema that already serves its child, so the admission is a download decision rather than a design one. & GEO GSE1\allowbreak 04878 an\allowbreak d GSE104\allowbreak 903; Add\allowbreak gene 127\allowbreak 546; whi\allowbreak ch acces\allowbreak sion hol\allowbreak ds which\allowbreak  library\allowbreak  is unre\allowbreak solved a\allowbreak nd one p\allowbreak age chec\allowbreak k settle\allowbreak s it \\
\addlinespace[5pt]
2 & \textbf{Vaishnav 2022 deferred training libraries} \org{S. cerevisiae}\newline Vaishnav ED, de Boer CG, Molinet J, Yassour M, Fan L, Adiconis X, Thompson DA, Levin JZ, Cubillos FA, Regev A. Nature 2022;603:455-463.\newline \href{https://doi.org/10.1038/s41586-022-04506-6}{doi:\allowbreak 10.\allowbreak 1038/\allowbreak s41586-\allowbreak 022-\allowbreak 04506-\allowbreak 6} & Regulatory DNA band, and the single lowest-effort large win in the wave: every piece of machinery it needs is already written and serving its subset. & Zenodo 4\allowbreak 436477 (\allowbreak complex\_\allowbreak media\_\allowbreak trainin\allowbreak g\_\allowbreak data\_\allowbreak Glu.\allowbreak txt 3.\allowbreak 59 GB, \allowbreak defined\_\allowbreak media\_\allowbreak trainin\allowbreak g\_\allowbreak data\_\allowbreak SC\_\allowbreak Ura.\allowbreak txt 2.\allowbreak 47 GB);\allowbreak  GEO GSE\allowbreak 163045 a\allowbreak nd GSE16\allowbreak 3866 \\
\addlinespace[5pt]
3 & \textbf{DREAM 2022 random promoter challenge (Rafi 2024)} \org{S. cerevisiae}\newline Rafi AM, et al. A community effort to optimize sequence-based deep learning models of gene regulation. Nat Biotechnol 2024.\newline \href{https://www.ncbi.nlm.nih.gov/geo/query/acc.cgi?acc=GSE254493}{ncbi.\allowbreak nlm.\allowbreak nih.\allowbreak gov/\allowbreak geo/\allowbreak query/\allowbreak acc.\allowbreak cgi?\allowbreak acc=\allowbreak GSE2544\allowbreak 93} & Regulatory DNA band: the benchmark slice, which is what makes a model comparison possible at all, and it carries its own eight-category test design. & GEO GSE2\allowbreak 54493; Z\allowbreak enodo 10\allowbreak 633252 (\allowbreak v2, abou\allowbreak t 4.\allowbreak 0 GB), \allowbreak supersed\allowbreak ing Zeno\allowbreak do 73953\allowbreak 97 \\
\addlinespace[5pt]
4 & \textbf{Kotopka and Smolke 2020 (artificial promoter libraries)} \org{S. cerevisiae}\newline Kotopka BJ, Smolke CD. Model-driven generation of artificial yeast promoters. Nat Commun 2020;11:2113.\newline \href{https://doi.org/10.1038/s41467-020-15977-4}{doi:\allowbreak 10.\allowbreak 1038/\allowbreak s41467-\allowbreak 020-\allowbreak 15977-\allowbreak 4} & Regulatory DNA band: the only million-scale promoter library outside the GPRA lineage, which is what makes it a check on it rather than more of it. & Zenodo 3\allowbreak 735426 (\allowbreak Promoter\allowbreak  Design \allowbreak Data\_\allowbreak 2020-\allowbreak 03.\allowbreak zip, ab\allowbreak out 1.\allowbreak 1 GB) \\
\addlinespace[5pt]
5 & \textbf{Zhu 2024 (Pol II MASTER, TSS initiation efficiency)} \org{S. cerevisiae}\newline Zhu C, Vvedenskaya IO, Sze SH, Nickels BE, Kaplan CD. Nat Struct Mol Biol 2024;31:190-202.\newline \href{https://doi.org/10.1038/s41594-023-01171-9}{doi:\allowbreak 10.\allowbreak 1038/\allowbreak s41594-\allowbreak 023-\allowbreak 01171-\allowbreak 9} & Regulatory DNA band: a different readout from reporter level, so it extends the axis rather than deepening it, and at a million records it is the third-largest row in the wave. & GEO GSE1\allowbreak 85290 (p\allowbreak rocessed\allowbreak ); BioPr\allowbreak oject PR\allowbreak JNA76662\allowbreak 4 (raw) \\
\addlinespace[5pt]
6 & \textbf{Cuperus 2017 (random 5' UTR library)} \org{S. cerevisiae}\newline Cuperus JT, Groves B, Kuchina A, Rosenberg AB, Jojic N, Fields S, Seelig G. Genome Res 2017;27:2015-2028.\newline \href{https://genome.cshlp.org/content/27/12/2015}{genome.\allowbreak cshlp.\allowbreak org/\allowbreak content\allowbreak /\allowbreak 27/\allowbreak 12/\allowbreak 2015} & Regulatory DNA band: the only large 5' UTR library, and the only one in the band whose readout is growth rather than fluorescence. & GEO GSE1\allowbreak 04252 \\
\addlinespace[5pt]
7 & \textbf{Gui 2021 (position effects with transcriptomes)} \org{S. cerevisiae}\newline Gui, et al. Transcriptome analysis in yeast reveals the externality of position effects. Mol Biol Evol 2021.\newline \href{https://www.biorxiv.org/content/10.1101/2020.04.02.021162}{biorxiv.\allowbreak org/\allowbreak content\allowbreak /\allowbreak 10.\allowbreak 1101/\allowbreak 2020.\allowbreak 04.\allowbreak 02.\allowbreak 021162} & Regulatory DNA band: the sequence variable is purely the integration coordinate, and this is the one row in that tier whose phenotype is a vector. & no acces\allowbreak sion fou\allowbreak nd; need\allowbreak s a brow\allowbreak ser chec\allowbreak k of the\allowbreak  publish\allowbreak ed suppl\allowbreak ement\newline {\scriptsize (browser)} \\
\addlinespace[5pt]
8 & \textbf{Keren 2016 (promoter panel x named gene, pooled fitness)} \org{S. cerevisiae}\newline Keren L, et al. Massively parallel interrogation of the effects of gene expression levels on fitness. Cell 2016;166:1282-1294.\newline \href{https://www.ncbi.nlm.nih.gov/geo/query/acc.cgi?acc=GSE83936}{ncbi.\allowbreak nlm.\allowbreak nih.\allowbreak gov/\allowbreak geo/\allowbreak query/\allowbreak acc.\allowbreak cgi?\allowbreak acc=\allowbreak GSE8393\allowbreak 6} & Regulatory DNA band: the only large tabulated panel that links expression level to fitness gene by gene, and the only one whose perturbation is a swap at a native locus. & GEO GSE8\allowbreak 3936; Ar\allowbreak rayExpre\allowbreak ss E-\allowbreak GEOD-\allowbreak 83936; \allowbreak BioProje\allowbreak ct PRJNA\allowbreak 329129 \\
\addlinespace[5pt]
9 & \textbf{Shalem 2015 (3' end FACS-seq library)} \org{S. cerevisiae}\newline Shalem O, et al. 3' end determinants. PLoS Genet 2015;11:e1005147.\newline \href{https://doi.org/10.1371/journal.pgen.1005147}{doi:\allowbreak 10.\allowbreak 1371/\allowbreak journal\allowbreak .\allowbreak pgen.\allowbreak 1005147} & Regulatory DNA band: the terminator end of the same axis, with the sequence and the value in one released file. & PLoS Gen\allowbreak et S3 Ta\allowbreak ble (XLS\allowbreak X), no a\allowbreak ccession\allowbreak ; needs \allowbreak a supple\allowbreak ment fet\allowbreak ch and a\allowbreak  sha256 \allowbreak pin\newline {\scriptsize (browser)} \\
\addlinespace[5pt]
10 & \textbf{Renganaath 2020 (native promoter variant MPRA)} \org{S. cerevisiae}\newline Renganaath K, Cheung R, Day L, Kosuri S, Kruglyak L, Albert FW. eLife 2020;9:e62669.\newline \href{https://doi.org/10.7554/eLife.62669}{doi:\allowbreak 10.\allowbreak 7554/\allowbreak eLife.\allowbreak 62669} & Regulatory DNA band: real alleles rather than sampled sequence, measured on the same readout, which is the bridge between the synthetic and the natural axes. & GEO GSE1\allowbreak 55944 (r\allowbreak aw plus \allowbreak barcode \allowbreak assignme\allowbreak nts); th\allowbreak e MPRA d\allowbreak esign an\allowbreak d proces\allowbreak sed tabl\allowbreak es are e\allowbreak Life Sup\allowbreak plementa\allowbreak ry Files \\
\addlinespace[5pt]
11 & \textbf{Arg-CCU tRNA deep mutational scan} \org{S. cerevisiae}\newline Deep mutational scan of a single-copy Arg-CCU tRNA over 14 substitutions seen across post-whole-genome-duplication Saccharomyces. PMC6193533.\newline \href{https://pmc.ncbi.nlm.nih.gov/articles/PMC6193533}{PMC61935\allowbreak 33} & Regulatory DNA band: point substitutions in a non-coding gene, which is the one DMS shape the protein aggregates cannot hold. & journal \allowbreak suppleme\allowbreak nt, no a\allowbreak ccession\allowbreak  located\allowbreak ; needs \allowbreak a supple\allowbreak ment fet\allowbreak ch\newline {\scriptsize (browser)} \\
\addlinespace[5pt]
12 & \textbf{Keren 2013 (native promoter panel across conditions)} \org{S. cerevisiae}\newline Keren L, et al. Mol Syst Biol 2013.\newline \href{https://spell.yeastgenome.org/}{spell.\allowbreak yeastge\allowbreak nome.\allowbreak org} & Regulatory DNA band: the only row in it with a real condition axis on native sequence. & not loca\allowbreak ted; rea\allowbreak ched onl\allowbreak y throug\allowbreak h the Zr\allowbreak imec and\allowbreak  Karollu\allowbreak s valida\allowbreak tion set\allowbreak s\newline {\scriptsize (browser)} \\
\addlinespace[5pt]
13 & \textbf{Lubliner 2015 (core promoter variant panel)} \org{S. cerevisiae}\newline Lubliner S, Keren L, Segal E, et al. Core promoter sequence in yeast is a major determinant of expression level. Genome Res 2015;25:1008-1017.\newline \href{https://genome.cshlp.org/content/25/7/1008}{genome.\allowbreak cshlp.\allowbreak org/\allowbreak content\allowbreak /\allowbreak 25/\allowbreak 7/\allowbreak 1008} & Regulatory DNA band: the core promoter is a distinct regulatory element from the upstream activating region every other row varies. & Genome R\allowbreak es Suppl\allowbreak emental \allowbreak Table 1;\allowbreak  needs a\allowbreak  supplem\allowbreak ent fetc\allowbreak h\newline {\scriptsize (browser)} \\
\addlinespace[5pt]
14 & \textbf{Sharon 2012 (designed promoter library)} \org{S. cerevisiae}\newline Sharon E, Kalma Y, Sharp A, et al., Segal E. Nat Biotechnol 2012;30:521-530.\newline \href{https://doi.org/10.1038/nbt.2205}{doi:\allowbreak 10.\allowbreak 1038/\allowbreak nbt.\allowbreak 2205} & Regulatory DNA band: the controlled counterpart the random libraries lack, and the platform's own origin. & Nat Biot\allowbreak echnol j\allowbreak ournal s\allowbreak upplemen\allowbreak t, no ac\allowbreak cession;\allowbreak  needs a\allowbreak  supplem\allowbreak ent fetc\allowbreak h\newline {\scriptsize (browser)} \\
\addlinespace[5pt]
15 & \textbf{Yamanishi 2013 (genomic terminatome)} \org{S. cerevisiae}\newline Yamanishi, et al. Characterization of heterologous protein expression using over 3,000 yeast genomic terminators. SGD reference S000206762.\newline \href{https://www.yeastgenome.org/reference/S000206762}{yeastgen\allowbreak ome.\allowbreak org/\allowbreak referen\allowbreak ce/\allowbreak S000206\allowbreak 762} & Regulatory DNA band: the largest native-terminator panel, which is the one element of the axis no library in the band randomizes. & primary \allowbreak table no\allowbreak t locate\allowbreak d; needs\allowbreak  the pap\allowbreak er ident\allowbreak ified, t\allowbreak hen a su\allowbreak pplement\allowbreak  fetch\newline {\scriptsize (browser)} \\
\addlinespace[5pt]
16 & \textbf{Dvir 2013 (5' UTR positions -10 to -1)} \org{S. cerevisiae}\newline Dvir S, Velten L, Sharon E, Zeevi D, Carey LB, Weinberger A, Segal E. Proc Natl Acad Sci USA 2013;110:E2792-E2801.\newline \href{https://doi.org/10.1073/pnas.1222534110}{doi:\allowbreak 10.\allowbreak 1073/\allowbreak pnas.\allowbreak 1222534\allowbreak 110} & Regulatory DNA band: a controlled ten-base window, which is the designed counterpart to the half-million random 5' UTRs above it. & PNAS jou\allowbreak rnal sup\allowbreak plement,\allowbreak  no acce\allowbreak ssion; n\allowbreak eeds a s\allowbreak upplemen\allowbreak t fetch\newline {\scriptsize (browser)} \\
\addlinespace[5pt]
17 & \textbf{Wu 2017 (genome-wide position effects, 1,044 loci)} \org{S. cerevisiae}\newline Wu XL, et al. Genome-wide landscape of position effects on heterogeneous gene expression in Saccharomyces cerevisiae. Biotechnol Biofuels 2017.\newline \href{https://www.ncbi.nlm.nih.gov/pmc/articles/PMC5516366/}{PMC55163\allowbreak 66} & Regulatory DNA band by its sequence variable, the integration coordinate, even though the readout is one reporter value. & Table S2\allowbreak , stated\allowbreak  availab\allowbreak le from \allowbreak the corr\allowbreak espondin\allowbreak g author\allowbreak  on reas\allowbreak onable r\allowbreak equest\newline {\scriptsize (browser)} \\
\addlinespace[5pt]
18 & \textbf{ProteinGym yeast subset} \org{S. cerevisiae}\newline ProteinGym: deep mutational scanning scores from 217 assays over 186 proteins, distributed through ProteinGymR and ProteinGymPy.\newline \href{https://proteingym.org}{proteing\allowbreak ym.\allowbreak org} & Regulatory DNA band as the aggregate route into the mutational-scan tier, which is the sequence-to-function axis in coding space rather than regulatory space. & ProteinG\allowbreak ymR and \allowbreak ProteinG\allowbreak ymPy pac\allowbreak kage loa\allowbreak ders; ar\allowbreak chived o\allowbreak n Zenodo\allowbreak , record\allowbreak  number \allowbreak not capt\allowbreak ured \\
\addlinespace[5pt]
19 & \textbf{Zhang 2016 (genome-wide GFP position scan)} \org{S. cerevisiae}\newline Zhang, et al. Genome-wide GFP position-effect scan using the yeast knockout collection. Chin J Biotech 2016.\newline \href{https://wsws.ijournals.cn/cjbcn/article/pdf/gc16070901}{wsws.\allowbreak ijourna\allowbreak ls.\allowbreak cn/\allowbreak cjbcn/\allowbreak article\allowbreak /\allowbreak pdf/\allowbreak gc16070\allowbreak 901} & Regulatory DNA band by its sequence variable, and the only FACS-read position-effect scan. & no acces\allowbreak sion loc\allowbreak ated; Ch\allowbreak inese-\allowbreak languag\allowbreak e journa\allowbreak l PDF, n\allowbreak eeds a b\allowbreak rowser f\allowbreak etch\newline {\scriptsize (browser)} \\
\addlinespace[5pt]
20 & \textbf{Kita 2017 (context-stratified yeast eQTLs)} \org{S. cerevisiae}\newline Kita R, et al. Proc Natl Acad Sci USA 2017;114.\newline \href{https://www.pnas.org/doi/10.1073/pnas.1717421114}{pnas.\allowbreak org/\allowbreak doi/\allowbreak 10.\allowbreak 1073/\allowbreak pnas.\allowbreak 1717421\allowbreak 114} & Regulatory DNA band: variants stratified by the regulatory element they sit in, which is the only released set with that stratification. & PNAS sup\allowbreak plementa\allowbreak ry data \allowbreak file pna\allowbreak s.\allowbreak 1717421\allowbreak 114.\allowbreak sd01; n\allowbreak eeds a s\allowbreak upplemen\allowbreak t fetch\newline {\scriptsize (browser)} \\
\addlinespace[5pt]
21 & \textbf{Chen (earlier position-effect study, 482 loci)} \org{S. cerevisiae}\newline An earlier Chen study characterizing position effects at 482 loci, reported inside Wu 2017; primary citation not located.\newline \href{https://www.ncbi.nlm.nih.gov/pmc/articles/PMC5516366/}{PMC55163\allowbreak 66} & Regulatory DNA band by its sequence variable, the integration coordinate. & not loca\allowbreak ted; the\allowbreak  primary\allowbreak  paper m\allowbreak ust be i\allowbreak dentifie\allowbreak d first\newline {\scriptsize (browser)} \\
\addlinespace[5pt]
22 & \textbf{Wei 2017 (synthetic terminator panel)} \org{S. cerevisiae}\newline Wei, et al. Yeast synthetic terminators: fine regulation of strength through linker sequences. ChemBioChem 2017.\newline \href{https://doi.org/10.1002/cbic.201700516}{doi:\allowbreak 10.\allowbreak 1002/\allowbreak cbic.\allowbreak 2017005\allowbreak 16} & Regulatory DNA band: the terminator parts tier, whose values are relative to a standard part rather than absolute. & ChemBioC\allowbreak hem jour\allowbreak nal supp\allowbreak lement, \allowbreak no repos\allowbreak itory de\allowbreak posit fo\allowbreak und\newline {\scriptsize (browser)} \\
\addlinespace[5pt]
23 & \textbf{Lubliner 2013 (native core promoter maximal strengths)} \org{S. cerevisiae}\newline Lubliner S, Keren L, Segal E. Nucleic Acids Res 2013.\newline \href{https://academic.oup.com/nar}{academic\allowbreak .\allowbreak oup.\allowbreak com/\allowbreak nar} & Regulatory DNA band: the native counterpart of the core-promoter tier, and the only source for its reference values. & not loca\allowbreak ted\newline {\scriptsize (browser)} \\
\addlinespace[5pt]
24 & \textbf{Morse (short synthetic terminators)} \org{S. cerevisiae}\newline Morse NJ, et al. Short synthetic terminators, 35 to 70 bp; primary citation not located.\newline \href{https://scite.ai/authors/nicholas-j-morse-1g5Qrn}{scite.\allowbreak ai/\allowbreak authors\allowbreak /\allowbreak nichola\allowbreak s-\allowbreak j-\allowbreak morse-\allowbreak 1g5Qrn} & Regulatory DNA band: the short-terminator corner of the parts tier, held so the lead is not lost. & no repos\allowbreak itory de\allowbreak posit fo\allowbreak und; pri\allowbreak mary pap\allowbreak er not i\allowbreak dentifie\allowbreak d\newline {\scriptsize (browser)} \\
\addlinespace[5pt]
25 & \textbf{MaveDB yeast score sets} \org{S. cerevisiae}\newline Esposito D, et al. MaveDB: an open-source platform to distribute and interpret data from multiplexed assays of variant effect. Genome Biol 2019;20:223.\newline \href{https://doi.org/10.1186/s13059-019-1845-6}{doi:\allowbreak 10.\allowbreak 1186/\allowbreak s13059-\allowbreak 019-\allowbreak 1845-\allowbreak 6} & Regulatory DNA band as the second aggregate route into the mutational-scan tier, held because the two aggregates do not hold the same assays. & mavedb.\allowbreak org; th\allowbreak e docume\allowbreak nted sea\allowbreak rch API \allowbreak endpoint\allowbreak  returns\allowbreak  HTTP 40\allowbreak 4, so th\allowbreak e organi\allowbreak sm filte\allowbreak r needs \allowbreak a live q\allowbreak uery\newline {\scriptsize (browser)} \\
\addlinespace[5pt]
26 & \textbf{Fitness Browser RCH2 (Stutzerimonas stutzeri\_AE)} \org{S. stutzeri}\newline Price MN, Wetmore KM, Waters RJ, Callaghan M, Ray J, Liu H, Kuehl JV, Melnyk RA, Lamson JS, Suh Y, et al., Deutschbauer AM, Arkin AP. Nature 2018;557:503-509.\newline \href{https://doi.org/10.1038/s41586-018-0124-0}{doi:\allowbreak 10.\allowbreak 1038/\allowbreak s41586-\allowbreak 018-\allowbreak 0124-\allowbreak 0} & Cross-host fitness: nine new hosts arrive inside one artifact under one schema, one condition vocabulary and an identifier bridge the source supplies, so the band exists to keep them together rather than scattered down a measurement ranking. & feba.\allowbreak db at 1\allowbreak 0,433,70\allowbreak 8,032 by\allowbreak tes, Las\allowbreak t-\allowbreak Modifie\allowbreak d 2026-\allowbreak 10-\allowbreak 27, HTT\allowbreak P 200 to\allowbreak  a plain\allowbreak  curl; o\allowbreak rgId psR\allowbreak CH2; ref\allowbreak seq\_\allowbreak mapping\allowbreak .\allowbreak tsv bri\allowbreak dges loc\allowbreak usId to \allowbreak RefSeq \\
\addlinespace[5pt]
27 & \textbf{Fitness Browser FW300-N2E3 (Aquipseudomonas helvetica)} \org{P. fluorescens}\newline Price MN, Wetmore KM, Waters RJ, et al., Deutschbauer AM, Arkin AP. Nature 2018;557:503-509.\newline \href{https://doi.org/10.1038/s41586-018-0124-0}{doi:\allowbreak 10.\allowbreak 1038/\allowbreak s41586-\allowbreak 018-\allowbreak 0124-\allowbreak 0} & Cross-host fitness: part of the single-artifact block of nine hosts, whose shared schema is the reason to admit them together. & feba.\allowbreak db, org\allowbreak Id pseud\allowbreak o3\_\allowbreak N2E3; t\allowbreak he histo\allowbreak rical bi\allowbreak nomial i\allowbreak s P.\allowbreak  fluore\allowbreak scens an\allowbreak d the cu\allowbreak rrent on\allowbreak e is A.\allowbreak  helvet\allowbreak ica \\
\addlinespace[5pt]
28 & \textbf{Fitness Browser FW300-N2E2 (Aquipseudomonas fluorescens\_Q)} \org{P. fluorescens}\newline Price MN, Wetmore KM, Waters RJ, et al., Deutschbauer AM, Arkin AP. Nature 2018;557:503-509.\newline \href{https://doi.org/10.1038/s41586-018-0124-0}{doi:\allowbreak 10.\allowbreak 1038/\allowbreak s41586-\allowbreak 018-\allowbreak 0124-\allowbreak 0} & Cross-host fitness: one of four isolates of one species in the block, which is the within-species comparison the band makes possible. & feba.\allowbreak db, org\allowbreak Id pseud\allowbreak o6\_\allowbreak N2E2; 1\allowbreak 16 updat\allowbreak ed annot\allowbreak ations \\
\addlinespace[5pt]
29 & \textbf{Fitness Browser WCS417 (Aquipseudomonas simiae)} \org{P. simiae}\newline Price MN, Wetmore KM, Waters RJ, et al., Deutschbauer AM, Arkin AP. Nature 2018;557:503-509.\newline \href{https://doi.org/10.1038/s41586-018-0124-0}{doi:\allowbreak 10.\allowbreak 1038/\allowbreak s41586-\allowbreak 018-\allowbreak 0124-\allowbreak 0} & Cross-host fitness: the plant-associated member of the single-artifact block, with its widest carbon panel. & feba.\allowbreak db, org\allowbreak Id WCS41\allowbreak 7; histo\allowbreak rical bi\allowbreak nomial P\allowbreak .\allowbreak  simiae\allowbreak , curren\allowbreak t A.\allowbreak  simiae \\
\addlinespace[5pt]
30 & \textbf{Fitness Browser FW300-N2C3 (Aquipseudomonas fluorescens\_AA)} \org{P. fluorescens}\newline Price MN, Wetmore KM, Waters RJ, et al., Deutschbauer AM, Arkin AP. Nature 2018;557:503-509.\newline \href{https://doi.org/10.1038/s41586-018-0124-0}{doi:\allowbreak 10.\allowbreak 1038/\allowbreak s41586-\allowbreak 018-\allowbreak 0124-\allowbreak 0} & Cross-host fitness: part of the four-isolate within-species block, and the row that taught the most unannotated genes. & feba.\allowbreak db, org\allowbreak Id pseud\allowbreak o5\_\allowbreak N2C3\_\allowbreak 1 \\
\addlinespace[5pt]
31 & \textbf{Fitness Browser FW300-N1B4 (Aquipseudomonas sp.)} \org{P. fluorescens}\newline Price MN, Wetmore KM, Waters RJ, et al., Deutschbauer AM, Arkin AP. Nature 2018;557:503-509.\newline \href{https://doi.org/10.1038/s41586-018-0124-0}{doi:\allowbreak 10.\allowbreak 1038/\allowbreak s41586-\allowbreak 018-\allowbreak 0124-\allowbreak 0} & Cross-host fitness: completes the four-isolate within-species block, and names the case of a host with no species assignment. & feba.\allowbreak db, org\allowbreak Id pseud\allowbreak o1\_\allowbreak N1B4 \\
\addlinespace[5pt]
32 & \textbf{Fitness Browser B728a (Aquipseudomonas fragariae pv. syringae)} \org{P. syringae}\newline Price MN, Wetmore KM, Waters RJ, et al., Deutschbauer AM, Arkin AP. Nature 2018;557:503-509.\newline \href{https://doi.org/10.1038/s41586-018-0124-0}{doi:\allowbreak 10.\allowbreak 1038/\allowbreak s41586-\allowbreak 018-\allowbreak 0124-\allowbreak 0} & Cross-host fitness: the pathogen member of the block, and half of the one strain in it screened in two genetic backgrounds. & feba.\allowbreak db, org\allowbreak Id Syrin\allowbreak gaeB728a\allowbreak ; the GT\allowbreak DB renam\allowbreak e to A.\allowbreak  fragar\allowbreak iae pv.\allowbreak  syring\allowbreak ae means\allowbreak  both bi\allowbreak nomials \allowbreak are need\allowbreak ed as al\allowbreak iases \\
\addlinespace[5pt]
33 & \textbf{Fitness Browser B728a mexB deletion background} \org{P. syringae}\newline Price MN, Wetmore KM, Waters RJ, et al., Deutschbauer AM, Arkin AP. Nature 2018;557:503-509.\newline \href{https://doi.org/10.1038/s41586-018-0124-0}{doi:\allowbreak 10.\allowbreak 1038/\allowbreak s41586-\allowbreak 018-\allowbreak 0124-\allowbreak 0} & Cross-host fitness: the background-crossed half of the B728a pair, which is the only interaction design in the block. & feba.\allowbreak db, org\allowbreak Id Syrin\allowbreak gaeB728a\_\allowbreak mexBdel\allowbreak ta \\
\addlinespace[5pt]
34 & \textbf{Fitness Browser PA14 (Pseudomonas aeruginosa)} \org{P. aeruginosa}\newline Price MN, Wetmore KM, Waters RJ, et al., Deutschbauer AM, Arkin AP. Nature 2018;557:503-509.\newline \href{https://doi.org/10.1038/s41586-018-0124-0}{doi:\allowbreak 10.\allowbreak 1038/\allowbreak s41586-\allowbreak 018-\allowbreak 0124-\allowbreak 0} & Cross-host fitness: the same artifact and schema, kept visible because its size is the finding about this species rather than an omission. & feba.\allowbreak db, org\allowbreak Id PA14;\allowbreak  the bro\allowbreak wser doe\allowbreak s not re\allowbreak name thi\allowbreak s organi\allowbreak sm \\
\addlinespace[5pt]
35 & \textbf{Dubinkina 2024 (12-strain low-pH transcriptomic atlas)} \org{P. kudriavzevii}\newline Dubinkina V, Bhogale S, Hsieh PH, Dibaeinia P, Nambiar S, Maslov S, Yoshikuni Y, Sinha S. Microbiol Spectr 2024.\newline \href{https://doi.org/10.1128/spectrum.02536-23}{doi:\allowbreak 10.\allowbreak 1128/\allowbreak spectru\allowbreak m.\allowbreak 02536-\allowbreak 23} & Host bootstrap: the only multi-strain expression matrix this host has, and the only deposit that already carries processed values rather than reads. & Illinois\allowbreak  Data Ba\allowbreak nk DOI 1\allowbreak 0.\allowbreak 13012/\allowbreak B2IDB-\allowbreak 1336177\_\allowbreak V1, hol\allowbreak ding str\allowbreak ain regu\allowbreak lators, \allowbreak strains \allowbreak and gene\allowbreak  sets, a\allowbreak nd gene \allowbreak expressi\allowbreak on count\allowbreak s; the p\allowbreak age retu\allowbreak rns 403 \allowbreak to every\allowbreak  automat\allowbreak ed clien\allowbreak t\newline {\scriptsize (browser)} \\
\addlinespace[5pt]
36 & \textbf{MAESTRO 2025 (170-isolate WGS, 57-stressor growth panel)} \org{P. kudriavzevii}\newline Hsieh PH, Sasaki Y, Yang CI, Wu ZY, Steindorff AS, Haridas S, Ke J, Fatma Z, Zhao Z, Opulente DA, Deng S, Hittinger CT, Grigoriev IV, Dien B, Zhao H, Li YP, Yoshikuni Y. bioRxiv 2025;2025.08.29.673034.\newline \href{https://doi.org/10.1101/2025.08.29.673034}{doi:\allowbreak 10.\allowbreak 1101/\allowbreak 2025.\allowbreak 08.\allowbreak 29.\allowbreak 673034} & Host bootstrap: this host has no genome-wide screen of any kind, so its largest resource is a 161-isolate panel that the measurement rule ranks below every pseudomonad screen. The band is what keeps it in the first tranche. & no acces\allowbreak sion ret\allowbreak rieved:\allowbreak  bioRxi\allowbreak v return\allowbreak s 403 an\allowbreak d Europe\allowbreak  PMC ful\allowbreak lTextXML\allowbreak  for PPR\allowbreak 1078985 \allowbreak returns \allowbreak 500; per\allowbreak -\allowbreak isolate\allowbreak  IO\_\allowbreak NNN Bio\allowbreak Projects\allowbreak  are inf\allowbreak erred, n\allowbreak ot confi\allowbreak rmed\newline {\scriptsize (ask)} \\
\addlinespace[5pt]
37 & \textbf{Fatma 2023 (landing-pad locus screen)} \org{P. kudriavzevii}\newline Fatma Z, Tan SI, Boob A, Zhao H. A landing pad system for multicopy gene integration in Issatchenkia orientalis. bioRxiv 2023;2023.05.21.541627.\newline \href{https://doi.org/10.1101/2023.05.21.541627}{doi:\allowbreak 10.\allowbreak 1101/\allowbreak 2023.\allowbreak 05.\allowbreak 21.\allowbreak 541627} & Host bootstrap: the host's only measured per-locus fitness, which no scale ranking would ever place in a first tranche at 188 records. & a 23 KB \allowbreak package \allowbreak at uofi.\allowbreak box.\allowbreak com (st\allowbreak rains, p\allowbreak lasmids,\allowbreak  guide R\allowbreak NAs, pri\allowbreak mers) pl\allowbreak us githu\allowbreak b.\allowbreak com/\allowbreak HuiminZ\allowbreak hao/\allowbreak Landing\allowbreak -\allowbreak pad-\allowbreak model; \allowbreak no DOI a\allowbreak nd no ar\allowbreak chival d\allowbreak eposit\newline {\scriptsize (ask)} \\
\addlinespace[5pt]
38 & \textbf{Douglass 2018 (32 sequenced isolates, 4-antifungal MIC)} \org{P. kudriavzevii}\newline Douglass AP, Offei B, Braun-Galleani S, Coughlan AY, Martos AAR, Ortiz-Merino RA, Byrne KP, Wolfe KH. PLoS Pathog 2018;14:e1007138.\newline \href{https://www.ncbi.nlm.nih.gov/pmc/articles/PMC6053246/}{PMC60532\allowbreak 46} & Host bootstrap: at 128 values it is the host's best public benchmark, which is the measure of how thin the host is and why a band rather than a score has to carry it. & SRA SRP1\allowbreak 39299 (3\allowbreak 2 strain\allowbreak s); GenB\allowbreak ank CP02\allowbreak 8773-\allowbreak CP02877\allowbreak 8 (CBS57\allowbreak 3 refere\allowbreak nce) and\allowbreak  CP02853\allowbreak 1-\allowbreak CP02853\allowbreak 5 (CBS51\allowbreak 47); MIC\allowbreak  table i\allowbreak n the op\allowbreak en-\allowbreak access \allowbreak article \allowbreak body \\
\addlinespace[5pt]
39 & \textbf{GSE282019 (Delhi NICU five-year outbreak)} \org{P. kudriavzevii}\newline Genomic analyses reveal high diversity and rapid evolution of Pichia kudriavzevii within a neonatal intensive care unit in Delhi, India. GEO series GSE282019, 2024; PMID 39853119.\newline \href{https://www.ncbi.nlm.nih.gov/geo/query/acc.cgi?acc=GSE282019}{ncbi.\allowbreak nlm.\allowbreak nih.\allowbreak gov/\allowbreak geo/\allowbreak query/\allowbreak acc.\allowbreak cgi?\allowbreak acc=\allowbreak GSE2820\allowbreak 19} & Host bootstrap: the host's only published drug-stratified expression series, and openly retrievable. & GEO GSE2\allowbreak 82019, p\allowbreak latform \allowbreak GPL34351\allowbreak ; progra\allowbreak mmatical\allowbreak ly retri\allowbreak evable \\
\addlinespace[5pt]
40 & \textbf{GSE202655 (L-malic acid, EcSthA transhydrogenase)} \org{P. kudriavzevii}\newline Metabolic engineering of the acid-tolerant yeast Pichia kudriavzevii for efficient L-malic acid production at low pH. GEO series GSE202655, 2023; PMID 36566973.\newline \href{https://www.ncbi.nlm.nih.gov/geo/query/acc.cgi?acc=GSE202655}{ncbi.\allowbreak nlm.\allowbreak nih.\allowbreak gov/\allowbreak geo/\allowbreak query/\allowbreak acc.\allowbreak cgi?\allowbreak acc=\allowbreak GSE2026\allowbreak 55} & Host bootstrap: an engineered-strain expression series in a non-SD108 background, which the host otherwise lacks. & GEO GSE2\allowbreak 02655, p\allowbreak latform \allowbreak GPL32243\allowbreak ; progra\allowbreak mmatical\allowbreak ly retri\allowbreak evable \\
\addlinespace[5pt]
41 & \textbf{GSE159457 (CY902 C4 dicarboxylic acid transporters)} \org{P. kudriavzevii}\newline Identification and characterization of C4 dicarboxylic acid transporters in Pichia kudriavzevii CY902. GEO series GSE159457, 2021; PMID 33629807.\newline \href{https://www.ncbi.nlm.nih.gov/geo/query/acc.cgi?acc=GSE159457}{ncbi.\allowbreak nlm.\allowbreak nih.\allowbreak gov/\allowbreak geo/\allowbreak query/\allowbreak acc.\allowbreak cgi?\allowbreak acc=\allowbreak GSE1594\allowbreak 57} & Host bootstrap: a carbon-source axis on expression in a third background, fully deposited. & GEO GSE1\allowbreak 59457, p\allowbreak latform \allowbreak GPL29244\allowbreak ; raw re\allowbreak ads at S\allowbreak RA SRP28\allowbreak 7255; pr\allowbreak ogrammat\allowbreak ically r\allowbreak etrievab\allowbreak le \\
\addlinespace[5pt]
42 & \textbf{GSE262836 (Delhi NICU, fluconazole-resistance expression)} \org{P. kudriavzevii}\newline Genomic characterization and transcriptomic profile of Candida krusei in a neonatal intensive care unit in Delhi, India. GEO series GSE262836, 2025.\newline \href{https://www.ncbi.nlm.nih.gov/geo/query/acc.cgi?acc=GSE262836}{ncbi.\allowbreak nlm.\allowbreak nih.\allowbreak gov/\allowbreak geo/\allowbreak query/\allowbreak acc.\allowbreak cgi?\allowbreak acc=\allowbreak GSE2628\allowbreak 36} & Host bootstrap: one of two rows in this host pairing a drug phenotype with expression, and openly retrievable. & GEO GSE2\allowbreak 62836, p\allowbreak latform \allowbreak GPL34351\allowbreak ; progra\allowbreak mmatical\allowbreak ly retri\allowbreak evable \\
\addlinespace[5pt]
43 & \textbf{Thompson 2019 (KT2440 lysine metabolism)} \org{P. putida}\newline Thompson MG, et al. Massively parallel fitness profiling reveals multiple novel enzymes in Pseudomonas putida lysine metabolism. mBio 2019.\newline \href{https://pmc.ncbi.nlm.nih.gov/articles/PMC6509195}{PMC65091\allowbreak 95} & Scale band: tier 3, ordered by measurements (22,256). & Table S1\allowbreak , a CSV \allowbreak suppleme\allowbreak ntary fi\allowbreak le; no r\allowbreak epositor\allowbreak y\newline {\scriptsize (browser)} \\
\addlinespace[5pt]
44 & \textbf{Isoprenol CRISPRi DBTL proteome ladder (7 PRIDE deposits)} \org{P. putida}\newline Automation of a CRISPRi platform for enhanced isoprenol production in Pseudomonas putida; PRIDE projects PXD063733, PXD063737, PXD063738, PXD063740, PXD063743, PXD063744, PXD063746.\newline \href{https://www.ebi.ac.uk/pride/archive/projects/PXD063733}{ebi.\allowbreak ac.\allowbreak uk/\allowbreak pride/\allowbreak archive\allowbreak /\allowbreak project\allowbreak s/\allowbreak PXD0637\allowbreak 33} & Scale band: tier 3, ordered by measurements (38,948). & PRIDE PX\allowbreak D063733,\allowbreak  PXD0637\allowbreak 37, PXD0\allowbreak 63738, P\allowbreak XD063740\allowbreak , PXD063\allowbreak 743, PXD\allowbreak 063744 a\allowbreak nd PXD06\allowbreak 3746, al\allowbreak l submit\allowbreak ted 2025\allowbreak -\allowbreak 05-\allowbreak 07 and \allowbreak 08 on an\allowbreak  Orbitra\allowbreak p Explor\allowbreak is 480 \\
\addlinespace[5pt]
45 & \textbf{Thompson 2020 (KT2440 fatty acids and alcohols)} \org{P. putida}\newline Thompson MG, Incha MR, Pearson AN, et al. Appl Environ Microbiol 2020;86:e01665-20; correction at 10.1128/AEM.00177-21.\newline \href{https://pubmed.ncbi.nlm.nih.gov/32826213}{PMID 328\allowbreak 26213} & Scale band: tier 3, ordered by measurements (127,972). & not stat\allowbreak ed; the \allowbreak values m\allowbreak ost like\allowbreak ly reach\allowbreak  feba.\allowbreak db \\
\addlinespace[5pt]
46 & \textbf{GSE198395 (putidaPRECISE321 carbon sources)} \org{P. putida}\newline GEO series GSE198395, KT2440 and engineered strains on different carbon sources, part of putidaPRECISE321, released 2022-04-28.\newline \href{https://www.ncbi.nlm.nih.gov/geo/query/acc.cgi?acc=GSE198395}{ncbi.\allowbreak nlm.\allowbreak nih.\allowbreak gov/\allowbreak geo/\allowbreak query/\allowbreak acc.\allowbreak cgi?\allowbreak acc=\allowbreak GSE1983\allowbreak 95} & Scale band: tier 3, ordered by measurements (450,684). & GEO GSE1\allowbreak 98395; p\allowbreak rogramma\allowbreak tically \allowbreak retrieva\allowbreak ble \\
\addlinespace[5pt]
47 & \textbf{Mueller 2025 (24-strain P. putida Biolog panel)} \org{P. putida}\newline Mueller J, Krishnan J, Wei Q, Hefner Y, Monk JM, Verkler H, Tibocha-Bonilla JD, et al. bioRxiv 2025;2025.11.17.688870.\newline \href{https://doi.org/10.1101/2025.11.17.688870}{doi:\allowbreak 10.\allowbreak 1101/\allowbreak 2025.\allowbreak 11.\allowbreak 17.\allowbreak 688870} & Scale band: tier 2, ordered by measurements (4,920). & unconfir\allowbreak med:\allowbreak  bioRxi\allowbreak v return\allowbreak ed HTTP \allowbreak 429 on f\allowbreak our atte\allowbreak mpts, so\allowbreak  no BioP\allowbreak roject, \allowbreak Zenodo r\allowbreak ecord or\allowbreak  machine\allowbreak -\allowbreak readabl\allowbreak e Table \allowbreak S2 is co\allowbreak nfirmed\newline {\scriptsize (ask)} \\
\addlinespace[5pt]
48 & \textbf{Schmidt 2022 (KT2440 nitrogen BarSeq)} \org{P. putida}\newline Schmidt M, Pearson AN, Incha MR, et al. Appl Environ Microbiol 2022;88:e02430-21.\newline \href{https://pmc.ncbi.nlm.nih.gov/articles/PMC9004399}{PMC90043\allowbreak 99} & Scale band: tier 2, ordered by measurements (395,044). & no acces\allowbreak sion sta\allowbreak ted; the\allowbreak  study l\allowbreak inks int\allowbreak o the Fi\allowbreak tness Br\allowbreak owser, s\allowbreak o its ex\allowbreak periment\allowbreak s are mo\allowbreak st likel\allowbreak y inside\allowbreak  feba.\allowbreak db \\
\addlinespace[5pt]
49 & \textbf{Borchert 2024 (KT2440 fitness compendium)} \org{P. putida}\newline Borchert AJ, Bleem AC, Lim HG, Rychel K, Dooley KD, Kellermyer ZA, Hodges TL, Palsson BO, Beckham GT. mSystems 2024;9:e00942-23.\newline \href{https://doi.org/10.1128/msystems.00942-23}{doi:\allowbreak 10.\allowbreak 1128/\allowbreak msystem\allowbreak s.\allowbreak 00942-\allowbreak 23} & Scale band: tier 1, ordered by measurements (1,571,024). & github.\allowbreak com/\allowbreak beckham\allowbreak -\allowbreak lab/\allowbreak fModule\allowbreak  (fitnes\allowbreak s values\allowbreak , statis\allowbreak tics, me\allowbreak tadata);\allowbreak  the non\allowbreak -\allowbreak browser\allowbreak  samples\allowbreak  are at \allowbreak SRA PRJN\allowbreak A809672,\allowbreak  PRJNA85\allowbreak 6070, PR\allowbreak JNA10112\allowbreak 87 \\
\addlinespace[5pt]
50 & \textbf{putidaPRECISE321 (Lim 2022)} \org{P. putida}\newline Lim HG, Rychel K, Sastry AV, Mueller J, Niu W, Feist AM, Palsson BO. Metab Eng 2022;72.\newline \href{https://doi.org/10.1016/j.ymben.2022.04.004}{doi:\allowbreak 10.\allowbreak 1016/\allowbreak j.\allowbreak ymben.\allowbreak 2022.\allowbreak 04.\allowbreak 004} & Scale band: tier 1, ordered by measurements (1,786,044). & github.\allowbreak com/\allowbreak SBRG/\allowbreak modulom\allowbreak e\_\allowbreak ppu (MI\allowbreak T); name\allowbreak d GEO de\allowbreak posits G\allowbreak SE198395\allowbreak  (81 sam\allowbreak ples) an\allowbreak d GSE193\allowbreak 493 (22)\allowbreak ; the fu\allowbreak ll per-\allowbreak sample \allowbreak GEO list\allowbreak  is in t\allowbreak he SI sp\allowbreak readshee\allowbreak t \\
\addlinespace[5pt]
51 & \textbf{Cao 2020 (pooled four-condition RNA-seq)} \org{P. kudriavzevii}\newline Cao M, et al., Zhao H. Metab Eng 2020;59:87-97, promoter-nomination RNA-seq.\newline \href{https://doi.org/10.1016/j.ymben.2020.01.005}{doi:\allowbreak 10.\allowbreak 1016/\allowbreak j.\allowbreak ymben.\allowbreak 2020.\allowbreak 01.\allowbreak 005} & Host bootstrap: the host's reference expression layer, and the one row that forces a decision about an environment that is a pool. & Illinois\allowbreak  Data Ba\allowbreak nk IDB-\allowbreak 8429399\allowbreak , Transc\allowbreak riptomic\allowbreak s tag; n\allowbreak o SRA or\allowbreak  GEO acc\allowbreak ession n\allowbreak amed; th\allowbreak e page r\allowbreak eturns 4\allowbreak 03\newline {\scriptsize (browser)} \\
\addlinespace[5pt]
52 & \textbf{PXD069498 (extracellular vesicle proteome)} \org{P. kudriavzevii}\newline PRIDE project PXD069498, submitter Chiara D'Ambrosio, April 2026.\newline \href{https://www.ebi.ac.uk/pride/archive/projects/PXD069498}{ebi.\allowbreak ac.\allowbreak uk/\allowbreak pride/\allowbreak archive\allowbreak /\allowbreak project\allowbreak s/\allowbreak PXD0694\allowbreak 98} & Host bootstrap: the host's only vesicle proteome, held for completeness at the bottom of the band. & PRIDE PX\allowbreak D069498;\allowbreak  retriev\allowbreak able thr\allowbreak ough the\allowbreak  PRIDE A\allowbreak PI \\
\addlinespace[5pt]
53 & \textbf{Douglass 2018 (per-strain SNP density and loss of heterozygosity)} \org{P. kudriavzevii}\newline Douglass AP, et al., Wolfe KH. PLoS Pathog 2018;14:e1007138, S2 Table.\newline \href{https://www.ncbi.nlm.nih.gov/pmc/articles/PMC6053246/}{PMC60532\allowbreak 46} & Host bootstrap: a genome-architecture layer for the host's only sequenced panel, which nothing else supplies. & S2 Table\allowbreak  of the \allowbreak open-\allowbreak access \allowbreak article \\
\addlinespace[5pt]
54 & \textbf{Tran 2023 (succinate end-to-end pipeline)} \org{P. kudriavzevii}\newline Tran VG, Mishra S, Bhagwat SS, Shafaei S, Shen Y, Allen JL, Crosly BA, Tan SI, Fatma Z, Rabinowitz JD, Guest JS, Singh V, Zhao H. Nat Commun 2023;14.\newline \href{https://pmc.ncbi.nlm.nih.gov/articles/PMC10547785/}{PMC10547\allowbreak 785} & Host bootstrap: the host's deepest named-genotype record, which is worth ingesting for the genotype side rather than for record count. & no seque\allowbreak nce or r\allowbreak aw-\allowbreak data ac\allowbreak cession;\allowbreak  all val\allowbreak ues in t\allowbreak he Suppl\allowbreak ementary\allowbreak  Informa\allowbreak tion and\allowbreak  the Sou\allowbreak rce Data\allowbreak  file, w\allowbreak ith plas\allowbreak mid sequ\allowbreak ences in\allowbreak  Supplem\allowbreak entary D\allowbreak ata 4\newline {\scriptsize (browser)} \\
\addlinespace[5pt]
55 & \textbf{Cao 2020 (constitutive promoter-GFP panel)} \org{P. kudriavzevii}\newline Cao M, Fatma Z, Song X, Hsieh PH, Tran VG, Lyon WL, Sayadi M, Shao Z, Yoshikuni Y, Zhao H. Metab Eng 2020;59:87-97.\newline \href{https://doi.org/10.1016/j.ymben.2020.01.005}{doi:\allowbreak 10.\allowbreak 1016/\allowbreak j.\allowbreak ymben.\allowbreak 2020.\allowbreak 01.\allowbreak 005} & Host bootstrap: the host's only parts panel, and the source of its gene-name to model-id map. & Illinois\allowbreak  Data Ba\allowbreak nk IDB-\allowbreak 8429399\allowbreak  (part s\allowbreak equences\allowbreak  plus a \allowbreak top-\allowbreak 10 prom\allowbreak oter lis\allowbreak t) and t\allowbreak he Elsev\allowbreak ier supp\allowbreak lement z\allowbreak ip; both\allowbreak  return \allowbreak 403 to a\allowbreak utomated\allowbreak  clients\newline {\scriptsize (browser)} \\
\addlinespace[5pt]
56 & \textbf{Xiao 2014 (SD108 reductive-TCA succinate)} \org{P. kudriavzevii}\newline Xiao H, et al., Zhao H. Microb Cell Fact 2014;13:121.\newline \href{https://pmc.ncbi.nlm.nih.gov/articles/PMC4244060}{PMC42440\allowbreak 60} & Host bootstrap: the strain-provenance root of every other engineering row in this host. & no acces\allowbreak sion nam\allowbreak ed; open\allowbreak -\allowbreak access \allowbreak article \\
\addlinespace[5pt]
57 & \textbf{Xylose utilization across four natural backgrounds} \org{P. kudriavzevii}\newline Cas9 metabolic-engineering work placing xylose reductase, xylitol dehydrogenase and xylulokinase from Scheffersomyces stipitis into four I. orientalis backgrounds; primary citation not located.\newline \href{https://core.ac.uk/works/286100561}{core.\allowbreak ac.\allowbreak uk/\allowbreak works/\allowbreak 2861005\allowbreak 61} & Host bootstrap: one construct in four backgrounds is the background-effect design, and it exists nowhere else in this wave. & not loca\allowbreak ted; onl\allowbreak y an agg\allowbreak regator \allowbreak record w\allowbreak as reach\allowbreak ed, so t\allowbreak he prima\allowbreak ry paper\allowbreak  must be\allowbreak  identif\allowbreak ied firs\allowbreak t\newline {\scriptsize (browser)} \\
\addlinespace[5pt]
58 & \textbf{Illinois Data Bank IDB-9063411 (high-yield 3-HP)} \org{P. kudriavzevii}\newline Data for high yield production of 3-hydroxypropionic acid using Issatchenkia orientalis. Illinois Data Bank, updated April 2026.\newline \href{https://databank.illinois.edu/datasets/IDB-9063411}{databank\allowbreak .\allowbreak illinoi\allowbreak s.\allowbreak edu/\allowbreak dataset\allowbreak s/\allowbreak IDB-\allowbreak 9063411} & Host bootstrap: a second product line on the host, with the highest yield it has recorded. & Illinois\allowbreak  Data Ba\allowbreak nk IDB-\allowbreak 9063411\allowbreak ; the pa\allowbreak ge retur\allowbreak ns 403\newline {\scriptsize (browser)} \\
\addlinespace[5pt]
59 & \textbf{Douglass 2018 (CBS573 RNA-seq)} \org{P. kudriavzevii}\newline Douglass AP, et al., Wolfe KH. PLoS Pathog 2018;14:e1007138; SRA SRP139056.\newline \href{https://www.ncbi.nlm.nih.gov/sra/?term=SRP139056}{ncbi.\allowbreak nlm.\allowbreak nih.\allowbreak gov/\allowbreak sra/\allowbreak ?\allowbreak term=\allowbreak SRP1390\allowbreak 56} & Host bootstrap: the only expression data on the host's clean reference, and the anchor for the namespace crosswalk. & SRA SRP1\allowbreak 39056 \\
\addlinespace[5pt]
60 & \textbf{MAESTRO 2025 (gene4260/gene4261 paralog-pair deletion)} \org{P. kudriavzevii}\newline Hsieh PH, et al., Yoshikuni Y. bioRxiv 2025;2025.08.29.673034, CRISPR validation arm.\newline \href{https://doi.org/10.1101/2025.08.29.673034}{doi:\allowbreak 10.\allowbreak 1101/\allowbreak 2025.\allowbreak 08.\allowbreak 29.\allowbreak 673034} & Host bootstrap: a measured gene perturbation in a host that has none at genome scale. & same unr\allowbreak esolved \allowbreak MAESTRO \allowbreak deposit;\allowbreak  needs t\allowbreak he brows\allowbreak er fetch\allowbreak  of the \allowbreak preprint\newline {\scriptsize (ask)} \\
\addlinespace[5pt]
61 & \textbf{Illinois Data Bank IDB-0538107 (succinate pipeline deposit)} \org{P. kudriavzevii}\newline Illinois Data Bank deposit, authors include Huimin Zhao; published March 2026.\newline \href{https://databank.illinois.edu/datasets/IDB-0538107}{databank\allowbreak .\allowbreak illinoi\allowbreak s.\allowbreak edu/\allowbreak dataset\allowbreak s/\allowbreak IDB-\allowbreak 0538107} & Host bootstrap: the host's deposited process data, which the journal pages do not carry. & Illinois\allowbreak  Data Ba\allowbreak nk IDB-\allowbreak 0538107\allowbreak ; the pa\allowbreak ge retur\allowbreak ns 403 t\allowbreak o automa\allowbreak ted clie\allowbreak nts\newline {\scriptsize (browser)} \\
\addlinespace[5pt]
62 & \textbf{CABBI inorganic carbon sources (low-pH succinate kinetics)} \org{P. kudriavzevii}\newline Data for influence of inorganic carbon sources on low-pH succinic acid production by Issatchenkia orientalis: process insights and kinetic analysis. CABBI dataset page.\newline \href{https://cabbi.bio/datasets/}{cabbi.\allowbreak bio/\allowbreak dataset\allowbreak s} & Host bootstrap: a process-condition axis on the host's main product. & CABBI da\allowbreak taset pa\allowbreak ge; the \allowbreak deposit \allowbreak target w\allowbreak as not r\allowbreak ead\newline {\scriptsize (browser)} \\
\addlinespace[5pt]
63 & \textbf{Xylose 3-HP from cellulosic hydrolysates (2026)} \org{P. kudriavzevii}\newline Xylose metabolic engineering of Issatchenkia orientalis for 3-hydroxypropionic acid production. Bioresour Technol 2026.\newline \href{https://doi.org/10.1016/j.biortech.2026.134783}{doi:\allowbreak 10.\allowbreak 1016/\allowbreak j.\allowbreak biortec\allowbreak h.\allowbreak 2026.\allowbreak 134783} & Host bootstrap: a feedstock axis on the host, which the glucose rows cannot supply. & companio\allowbreak n CABBI \allowbreak dataset \allowbreak page; de\allowbreak posit ta\allowbreak rget not\allowbreak  read; n\allowbreak o PMCID\newline {\scriptsize (browser)} \\
\addlinespace[5pt]
64 & \textbf{Jeong 2025 (acetate metabolism during xylose fermentation)} \org{P. kudriavzevii}\newline Jeong D, et al. Acetate metabolism during xylose fermentation enhances 3-HP production in acid-tolerant I. orientalis. bioRxiv 2025;2025.04.17.649204.\newline \href{https://doi.org/10.1101/2025.04.17.649204}{doi:\allowbreak 10.\allowbreak 1101/\allowbreak 2025.\allowbreak 04.\allowbreak 17.\allowbreak 649204} & Host bootstrap: a third carbon arm on the host's second product. & bioRxiv \allowbreak only; ar\allowbreak ticle pa\allowbreak ges retu\allowbreak rn 403 t\allowbreak o automa\allowbreak ted fetc\allowbreak h\newline {\scriptsize (browser)} \\
\addlinespace[5pt]
65 & \textbf{Glycolic acid overproduction (2026)} \org{P. kudriavzevii}\newline Xylose-assisted glycolic acid overproduction in engineered Issatchenkia orientalis. J Fungi 2026;12:751.\newline \href{https://doi.org/10.3390/jof12100751}{doi:\allowbreak 10.\allowbreak 3390/\allowbreak jof1210\allowbreak 0751} & Host bootstrap: a fourth product class on the host, and one of the few rows in it that is openly retrievable. & open-\allowbreak access \allowbreak article \allowbreak (PMC1370\allowbreak 5920); n\allowbreak o separa\allowbreak te depos\allowbreak it named \\
\addlinespace[5pt]
66 & \textbf{Wu 2023 (piggyBac citramalate copy-number and locus variants)} \org{P. kudriavzevii}\newline Wu, et al. Metab Eng Commun 2023;16:e00220.\newline \href{https://pmc.ncbi.nlm.nih.gov/articles/PMC9969067}{PMC99690\allowbreak 67} & Host bootstrap: a copy-number by locus dose axis, which is the only quantitative genotype axis the host offers. & Illinois\allowbreak  Data Ba\allowbreak nk IDB-\allowbreak 6723545\allowbreak  (publis\allowbreak hed Octo\allowbreak ber 2025\allowbreak , contac\allowbreak t Yoshik\allowbreak uni); th\allowbreak e page r\allowbreak eturns 4\allowbreak 03; the \allowbreak article \allowbreak is open \allowbreak access\newline {\scriptsize (browser)} \\
\addlinespace[5pt]
67 & \textbf{Mitochondrial decompartmentalization (2025, with GNPS metabolomics)} \org{P. kudriavzevii}\newline Decompartmentalization of the yeast mitochondrial metabolism to improve chemical production in Issatchenkia orientalis. Nat Commun 2025.\newline \href{https://doi.org/10.1038/s41467-025-62304-w}{doi:\allowbreak 10.\allowbreak 1038/\allowbreak s41467-\allowbreak 025-\allowbreak 62304-\allowbreak w} & Host bootstrap: the host's only metabolomics layer. & a GNPS d\allowbreak epositio\allowbreak n named \allowbreak on the C\allowbreak ABBI dat\allowbreak aset pag\allowbreak e but no\allowbreak t read; \allowbreak a furthe\allowbreak r Illino\allowbreak is Data \allowbreak Bank dep\allowbreak osit sur\allowbreak faces on\allowbreak ly throu\allowbreak gh OSTI \allowbreak as DOI 1\allowbreak 0.\allowbreak 13012/\allowbreak B2IDB-\allowbreak 9575028\_\allowbreak V1\newline {\scriptsize (browser)} \\
\addlinespace[5pt]
68 & \textbf{GSE39558 (ethanol-stress expression)} \org{P. kudriavzevii}\newline GEO series GSE39558, ethanol-stress RNA-seq; species assignment reported by a search index only.\newline \href{https://www.ncbi.nlm.nih.gov/geo/query/acc.cgi?acc=GSE39558}{ncbi.\allowbreak nlm.\allowbreak nih.\allowbreak gov/\allowbreak geo/\allowbreak query/\allowbreak acc.\allowbreak cgi?\allowbreak acc=\allowbreak GSE3955\allowbreak 8} & Host bootstrap: a stress axis the host does not otherwise have in expression, and the cheapest row in the wave to verify. & GEO GSE3\allowbreak 9558; on\allowbreak e call v\allowbreak erifies \allowbreak contents\allowbreak , sample\allowbreak  count a\allowbreak nd strai\allowbreak n \\
\addlinespace[5pt]
69 & \textbf{PXD041801 (flocculin expression and floating biofilms)} \org{P. kudriavzevii}\newline Raw proteomics for a 2023 PLoS Pathogens study of flocculin expression and floating biofilms; submitted via MassIVE.\newline \href{https://www.ebi.ac.uk/pride/archive/projects/PXD041801}{ebi.\allowbreak ac.\allowbreak uk/\allowbreak pride/\allowbreak archive\allowbreak /\allowbreak project\allowbreak s/\allowbreak PXD0418\allowbreak 01} & Host bootstrap: a biofilm axis the host has nowhere else. & Proteome\allowbreak Xchange \allowbreak PXD04180\allowbreak 1 via Ma\allowbreak ssIVE; r\allowbreak etrievab\allowbreak le \\
\addlinespace[5pt]
70 & \textbf{PXD039973 (succinylation and amphotericin B resistance)} \org{P. kudriavzevii}\newline PRIDE/iProX project PXD039973, protein succinylation linked to amphotericin B resistance.\newline \href{https://www.omicsdi.org/dataset/iprox/PXD039973}{omicsdi.\allowbreak org/\allowbreak dataset\allowbreak /\allowbreak iprox/\allowbreak PXD0399\allowbreak 73} & Host bootstrap: the host's only modification-site layer, and the only one tied to a drug phenotype. & PRIDE/\allowbreak iProX P\allowbreak XD039973\allowbreak ; iProX \allowbreak retrieva\allowbreak bility u\allowbreak nverifie\allowbreak d \\
\addlinespace[5pt]
71 & \textbf{PXD075355 (Pichia against Botrytis cinerea)} \org{P. kudriavzevii}\newline PRIDE project PXD075355, native Pichia against Botrytis cinerea, 2026; lists both P. kudriavzevii and P. kluyveri.\newline \href{https://www.ebi.ac.uk/pride/archive/projects/PXD075355}{ebi.\allowbreak ac.\allowbreak uk/\allowbreak pride/\allowbreak archive\allowbreak /\allowbreak project\allowbreak s/\allowbreak PXD0753\allowbreak 55} & Host bootstrap: the only biotic-interaction environment in the whole wave. & PRIDE PX\allowbreak D075355;\allowbreak  retriev\allowbreak able, bu\allowbreak t the de\allowbreak posit ho\allowbreak lds two \allowbreak species \\
\addlinespace[5pt]
72 & \textbf{Pk E1 toolkit expansion (mycoprotein, 2026)} \org{P. kudriavzevii}\newline Expanding the genetic toolkit of Pichia kudriavzevii E1 for enhanced mycoprotein production. Microb Biotechnol 2026;19.\newline \href{https://doi.org/10.1111/1751-7915.70447}{doi:\allowbreak 10.\allowbreak 1111/\allowbreak 1751-\allowbreak 7915.\allowbreak 70447} & Host bootstrap: a second parts panel in a second background, which is what tests whether the first one transfers. & open-\allowbreak access \allowbreak article \allowbreak (PMC1357\allowbreak 8925); n\allowbreak o separa\allowbreak te depos\allowbreak it named \\
\addlinespace[5pt]
73 & \textbf{China fermented-vegetable and clinical panel (2024)} \org{P. kudriavzevii}\newline Comparative genomics and susceptibility of fermented-vegetable and clinical Candida krusei isolates. Virulence 2024;15.\newline \href{https://www.ncbi.nlm.nih.gov/pmc/articles/PMC11487970/}{PMC11487\allowbreak 970} & Host bootstrap: the host's only second clinical panel, held because one supplement check decides whether it is a dataset. & open-\allowbreak access \allowbreak article;\allowbreak  per-\allowbreak isolate\allowbreak  MIC tab\allowbreak les and \allowbreak genome a\allowbreak ccession\allowbreak s not ve\allowbreak rified\newline {\scriptsize (browser)} \\
\addlinespace[5pt]
74 & \textbf{Indonesian clinical genomes (PRJNA819536)} \org{P. kudriavzevii}\newline Genome note on whole-genome sequencing of 16 clinical isolates from Jakarta hospitals. F1000Research 2022;11:593.\newline \href{https://www.ncbi.nlm.nih.gov/pmc/articles/PMC10320322/}{PMC10320\allowbreak 322} & Host bootstrap: it would add half again as many sequenced isolates to the host's only genotyped panel. & BioProje\allowbreak ct PRJNA\allowbreak 819536 f\allowbreak or the s\allowbreak equence;\allowbreak  the phe\allowbreak notype t\allowbreak able's l\allowbreak ocation \allowbreak is unkno\allowbreak wn and m\allowbreak ay be an\allowbreak  ask\newline {\scriptsize (browser)} \\
\addlinespace[5pt]
75 & \textbf{Unidentified Illinois Data Bank CRISPR/Cas9 gene-deletion deposit} \org{P. kudriavzevii}\newline An Illinois Data Bank CRISPR/Cas9 gene-deletion dataset surfacing only through an OSTI bibliographic record; deposit identifier not confirmed.\newline \href{https://www.osti.gov/biblio/3013685}{osti.\allowbreak gov/\allowbreak biblio/\allowbreak 3013685} & Host bootstrap: the one lead that could overturn the finding that this host has no multi-gene deletion resource. & an Illin\allowbreak ois Data\allowbreak  Bank de\allowbreak posit wh\allowbreak ose iden\allowbreak tifier i\allowbreak s unconf\allowbreak irmed; n\allowbreak eeds an \allowbreak OSTI and\allowbreak  a repos\allowbreak itory br\allowbreak owser vi\allowbreak sit\newline {\scriptsize (browser)} \\
\addlinespace[5pt]
76 & \textbf{Khaledi 2020 (414-isolate clinical transcriptome panel)} \org{P. aeruginosa}\newline Khaledi A, Weimann A, Schniederjans M, Asgari E, Kuo TH, Oliver A, Cabot G, Kola A, Gastmeier P, Hogardt M, Jonas D, Mofrad MRK, Bremges A, McHardy AC, Haussler S. EMBO Mol Med 2020;12:e10264.\newline \href{https://doi.org/10.15252/emmm.201910264}{doi:\allowbreak 10.\allowbreak 15252/\allowbreak emmm.\allowbreak 2019102\allowbreak 64} & Scale band: tier 1, ordered by measurements (2,442,600). & GEO GSE1\allowbreak 23544 (R\allowbreak NA-\allowbreak seq); B\allowbreak ioProjec\allowbreak t PRJNA5\allowbreak 26797 (D\allowbreak NA-\allowbreak seq); Z\allowbreak enodo 10\allowbreak .\allowbreak 5281/\allowbreak zenodo.\allowbreak 3464542\allowbreak  (classi\allowbreak fier inp\allowbreak uts); re\allowbreak ference \allowbreak NC\_\allowbreak 008463.\allowbreak 1 \\
\addlinespace[5pt]
77 & \textbf{Kaczmarczyk 2025 (PAO1 genome-wide CRISPRi)} \org{P. aeruginosa}\newline Kaczmarczyk A, Klotz A, Manfredi P, Jenal U. bioRxiv 2025;2025.08.12.669819.\newline \href{https://doi.org/10.1101/2025.08.12.669819}{doi:\allowbreak 10.\allowbreak 1101/\allowbreak 2025.\allowbreak 08.\allowbreak 12.\allowbreak 669819} & Scale band: tier 2, ordered by measurements (80,000). & Zenodo 1\allowbreak 0.\allowbreak 5281/\allowbreak zenodo.\allowbreak 1559174\allowbreak 2, .\allowbreak 1559223\allowbreak 8 and .\allowbreak 1559336\allowbreak 1, conte\allowbreak nts unla\allowbreak beled; n\allowbreak o SRA or\allowbreak  GEO acc\allowbreak ession\newline {\scriptsize (ask)} \\
\addlinespace[5pt]
78 & \textbf{Dotsch 2009 (arrayed PA14 library against 19 antibiotics)} \org{P. aeruginosa}\newline Dotsch A, Becker T, Pommerenke C, Janowska Z, Jansch L, Haussler S. Antimicrob Agents Chemother 2009;53:2522-2531.\newline \href{https://doi.org/10.1128/AAC.00035-09}{doi:\allowbreak 10.\allowbreak 1128/\allowbreak AAC.\allowbreak 00035-\allowbreak 09} & Scale band: tier 2, ordered by measurements (5,459). & value lo\allowbreak cation u\allowbreak nconfirm\allowbreak ed; the \allowbreak PA14 non\allowbreak -\allowbreak redunda\allowbreak nt libra\allowbreak ry backb\allowbreak one is a\allowbreak t the PA\allowbreak TIMDB po\allowbreak rtal, wh\allowbreak ose down\allowbreak loads pa\allowbreak ge was n\allowbreak ot read\newline {\scriptsize (ask)} \\
\addlinespace[5pt]
79 & \textbf{Bartell 2019 (longitudinal cystic fibrosis panel)} \org{P. aeruginosa}\newline Bartell JA, Sommer LM, Haagensen JAJ, Loch A, Espinosa R, Molin S, Johansen HK. Nat Commun 2019;10:629.\newline \href{https://doi.org/10.1038/s41467-019-08504-7}{doi:\allowbreak 10.\allowbreak 1038/\allowbreak s41467-\allowbreak 019-\allowbreak 08504-\allowbreak 7} & Scale band: tier 2, ordered by measurements (3,544). & Suppleme\allowbreak ntary Da\allowbreak ta 1 (xl\allowbreak sx) hold\allowbreak s the en\allowbreak tire dat\allowbreak aset and\allowbreak  the iso\allowbreak late acc\allowbreak essions;\allowbreak  genomes\allowbreak  at SRA \allowbreak study ER\allowbreak P004853\newline {\scriptsize (browser)} \\
\addlinespace[5pt]
80 & \textbf{Genome Medicine 2024 (1,877-strain genome and MIC set)} \org{P. aeruginosa}\newline Paired whole-genome and susceptibility data for 1,877 diverse strains plus a 102-strain validation set, with MICs reclassified per CLSI M100S-Ed32:2022. Genome Med 2024;16.\newline \href{https://doi.org/10.1186/s13073-024-01346-z}{doi:\allowbreak 10.\allowbreak 1186/\allowbreak s13073-\allowbreak 024-\allowbreak 01346-\allowbreak z} & Scale band: tier 2, ordered by measurements (1,979). & unconfir\allowbreak med; the\allowbreak  publish\allowbreak er redir\allowbreak ects to \allowbreak an authe\allowbreak nticatio\allowbreak n endpoi\allowbreak nt\newline {\scriptsize (browser)} \\
\addlinespace[5pt]
81 & \textbf{Penaranda 2025 (125-isolate diversity survey)} \org{P. aeruginosa}\newline Penaranda C, Brenner EP, Clatworthy AE, Cosimi LA, Ravi J, Hung DT. bioRxiv 2025;2025.09.22.677895.\newline \href{https://doi.org/10.1101/2025.09.22.677895}{doi:\allowbreak 10.\allowbreak 1101/\allowbreak 2025.\allowbreak 09.\allowbreak 22.\allowbreak 677895} & Scale band: tier 2, ordered by measurements (1,664). & Suppleme\allowbreak ntal Dat\allowbreak a 1 (xls\allowbreak x) holds\allowbreak  genome \allowbreak accessio\allowbreak ns, meta\allowbreak data and\allowbreak  phenoty\allowbreak pes; the\allowbreak  BioProj\allowbreak ect numb\allowbreak er is in\allowbreak side tha\allowbreak t file a\allowbreak nd uncon\allowbreak firmed\newline {\scriptsize (browser)} \\
\addlinespace[5pt]
82 & \textbf{Khaledi 2020 (1,656-record agar-dilution MIC panel)} \org{P. aeruginosa}\newline Khaledi A, et al., Haussler S. EMBO Mol Med 2020;12:e10264, susceptibility layer.\newline \href{https://doi.org/10.15252/emmm.201910264}{doi:\allowbreak 10.\allowbreak 15252/\allowbreak emmm.\allowbreak 2019102\allowbreak 64} & Scale band: tier 2, ordered by measurements (1,656). & phenotyp\allowbreak e table \allowbreak as Datas\allowbreak et EV1; \allowbreak Zenodo 1\allowbreak 0.\allowbreak 5281/\allowbreak zenodo.\allowbreak 3464542\allowbreak  for the\allowbreak  classif\allowbreak ier inpu\allowbreak ts \\
\addlinespace[5pt]
83 & \textbf{Poulsen 2019 (nine-strain Tn-seq core essentiality)} \org{P. aeruginosa}\newline Poulsen BE, et al. Defining the core essential genome of Pseudomonas aeruginosa. bioRxiv 2019;396689.\newline \href{https://www.biorxiv.org/content/10.1101/396689v2.full.pdf}{biorxiv.\allowbreak org/\allowbreak content\allowbreak /\allowbreak 10.\allowbreak 1101/\allowbreak 396689v\allowbreak 2.\allowbreak full.\allowbreak pdf} & Scale band: tier 2, ordered by measurements (321). & not give\allowbreak n by the\allowbreak  compili\allowbreak ng paper\allowbreak ; the pr\allowbreak eprint i\allowbreak s open\newline {\scriptsize (browser)} \\
\addlinespace[5pt]
84 & \textbf{PAO1 arrayed transposon libraries (Jacobs 2003, Held 2012)} \org{P. aeruginosa}\newline Jacobs MA, et al. Proc Natl Acad Sci USA 2003;100:14339-14344; Held K, Ramage E, Jacobs M, Gallagher L, Manoil C. J Bacteriol 2012;194.\newline \href{https://doi.org/10.1128/JB.01479-12}{doi:\allowbreak 10.\allowbreak 1128/\allowbreak JB.\allowbreak 01479-\allowbreak 12} & Scale band: tier 3, ordered by measurements (9,437). & insertio\allowbreak n-\allowbreak positio\allowbreak n metada\allowbreak ta from \allowbreak the libr\allowbreak ary inde\allowbreak x page; \allowbreak strains \allowbreak are orde\allowbreak r-\allowbreak only th\allowbreak rough a \allowbreak cost cen\allowbreak ter and \allowbreak distribu\allowbreak tion has\allowbreak  moved l\allowbreak abs\newline {\scriptsize (browser)} \\
\addlinespace[5pt]
85 & \textbf{KT2440 pooled transposon library (Putida\_ML5 metadata)} \org{P. putida}\newline Pooled RB-TnSeq library metadata for Pseudomonas putida KT2440, one mapped version named Putida\_ML5.\newline \href{https://nrel.gov/docs/fy22osti/83183.pdf}{nrel.\allowbreak gov/\allowbreak docs/\allowbreak fy22ost\allowbreak i/\allowbreak 83183.\allowbreak pdf} & Scale band: tier 3, ordered by measurements (4,778). & NREL met\allowbreak hods doc\allowbreak umentati\allowbreak on plus \allowbreak the mapp\allowbreak ing arch\allowbreak ives shi\allowbreak pped wit\allowbreak h the bi\allowbreak gfit rel\allowbreak ease \\
\addlinespace[5pt]
86 & \textbf{Gao 2020 (KT2440 targeted SRM proteomics)} \org{P. putida}\newline Gao Y, Fillmore TL, Munoz N, Bentley GJ, Johnson CW, et al. Front Bioeng Biotechnol 2020;8:603488.\newline \href{https://doi.org/10.3389/fbioe.2020.603488}{doi:\allowbreak 10.\allowbreak 3389/\allowbreak fbioe.\allowbreak 2020.\allowbreak 603488} & Scale band: tier 3, ordered by measurements (1,056). & Panorama\allowbreak  Public,\allowbreak  no PXD \allowbreak accessio\allowbreak n; monit\allowbreak ored pep\allowbreak tides in\allowbreak  Supplem\allowbreak entary T\allowbreak ables S1\allowbreak  and S2,\allowbreak  limits \allowbreak of quant\allowbreak ificatio\allowbreak n in S3 \\
\addlinespace[5pt]
87 & \textbf{BACTOME (101 clinical genomes with eight phenotypes)} \org{P. aeruginosa}\newline Hornischer K, et al. BACTOME: a reference database to explore the sequence- and gene expression-variation landscape of Pseudomonas aeruginosa clinical isolates. Nucleic Acids Res 2019;47:D716.\newline \href{https://doi.org/10.1093/nar/gky895}{doi:\allowbreak 10.\allowbreak 1093/\allowbreak nar/\allowbreak gky895} & Scale band: tier 3, ordered by measurements (808). & no bulk \allowbreak download\allowbreak  and no \allowbreak ENA, SRA\allowbreak  or GEO \allowbreak accessio\allowbreak n on the\allowbreak  landing\allowbreak  page; n\allowbreak ot confi\allowbreak rmed bul\allowbreak k-\allowbreak retriev\allowbreak able\newline {\scriptsize (ask)} \\
\addlinespace[5pt]
88 & \textbf{Kohlstedt 2019 (13C metabolic flux analysis, two hosts)} \org{P. putida}\newline Kohlstedt M, Wittmann C. Metab Eng 2019;54:35-53.\newline \href{https://doi.org/10.1016/j.ymben.2019.01.008}{doi:\allowbreak 10.\allowbreak 1016/\allowbreak j.\allowbreak ymben.\allowbreak 2019.\allowbreak 01.\allowbreak 008} & Scale band: tier 3, ordered by measurements (534). & unconfir\allowbreak med:\allowbreak  the pu\allowbreak blisher \allowbreak returned\allowbreak  403 and\allowbreak  the ins\allowbreak titution\allowbreak al PDF r\allowbreak eset the\allowbreak  connect\allowbreak ion\newline {\scriptsize (browser)} \\
\addlinespace[5pt]
89 & \textbf{Turner 2015 (PA14 and PAO1 essentiality across four arms)} \org{P. aeruginosa}\newline Turner KH, Wessel AK, Palmer GC, Murray JL, Whiteley M. Proc Natl Acad Sci USA 2015.\newline \href{https://pmc.ncbi.nlm.nih.gov/articles/PMC4386324}{PMC43863\allowbreak 24} & Scale band: tier 3, ordered by measurements (352). & SRA PRJN\allowbreak A265367 \\
\addlinespace[5pt]
90 & \textbf{Liberati 2006 (PA14 non-redundant transposon library and essentiality)} \org{P. aeruginosa}\newline Liberati NT, Urbach JM, Miyata S, Lee DG, Drenkard E, Wu G, Villanueva J, Wei T, Ausubel FM. Proc Natl Acad Sci USA 2006;103:2833-2838.\newline \href{https://doi.org/10.1073/pnas.0511100103}{doi:\allowbreak 10.\allowbreak 1073/\allowbreak pnas.\allowbreak 0511100\allowbreak 103} & Scale band: tier 3, ordered by measurements (335). & the PATI\allowbreak MDB port\allowbreak al, with\allowbreak  Search,\allowbreak  BLAST a\allowbreak nd Inser\allowbreak tion Sit\allowbreak e Map pa\allowbreak ges; whe\allowbreak ther per\allowbreak -\allowbreak mutant \allowbreak tables a\allowbreak re bulk-\allowbreak downloa\allowbreak dable is\allowbreak  unconfi\allowbreak rmed\newline {\scriptsize (browser)} \\
\addlinespace[5pt]
91 & \textbf{GOLD\_115 and GOLD\_84 essential-gene sets} \org{P. aeruginosa}\newline Van Maele C, Caboche S, Nicolau-Guillaumet N, Muggeo A, Guillard T. PLoS Comput Biol 2026;22:e1013945.\newline \href{https://doi.org/10.1371/journal.pcbi.1013945}{doi:\allowbreak 10.\allowbreak 1371/\allowbreak journal\allowbreak .\allowbreak pcbi.\allowbreak 1013945} & Scale band: tier 3, ordered by measurements (199). & S2 Table\allowbreak , Gold S\allowbreak tandard \allowbreak List; th\allowbreak e paper'\allowbreak s own re\allowbreak ads at S\allowbreak RA PRJNA\allowbreak 1240204;\allowbreak  referen\allowbreak ce NC\_\allowbreak 008463.\allowbreak 1\newline {\scriptsize (browser)} \\
\addlinespace[5pt]
92 & \textbf{Werner 2023 (beta-ketoadipate from lignin)} \org{P. putida}\newline Werner AZ, Cordell WT, Lahive CW, Klein BC, Singer CA, Tan ECD, et al., Beckham GT. Sci Adv 2023;9:eadj0053.\newline \href{https://doi.org/10.1126/sciadv.adj0053}{doi:\allowbreak 10.\allowbreak 1126/\allowbreak sciadv.\allowbreak adj0053} & Scale band: tier 3, ordered by measurements (90). & sciadv.\allowbreak adj0053\_\allowbreak data\_\allowbreak s1.\allowbreak zip plu\allowbreak s the SI\allowbreak  PDF; no\allowbreak  reposit\allowbreak ory; pla\allowbreak smids de\allowbreak ferred t\allowbreak o Tables\allowbreak  S1 to S\allowbreak 3\newline {\scriptsize (browser)} \\
\addlinespace[5pt]
93 & \textbf{Wilkes 2026 (muconate multi-omics with fluxomics)} \org{P. putida}\newline Wilkes RA, Suthers PF, Borchert AJ, Callaghan MM, Thusoo E, Giannone RJ, et al. Metab Eng 2026;99:102573.\newline \href{https://doi.org/10.1016/j.ymben.2026.102573}{doi:\allowbreak 10.\allowbreak 1016/\allowbreak j.\allowbreak ymben.\allowbreak 2026.\allowbreak 102573} & Scale band: tier 3, ordered by measurements (12). & none con\allowbreak firmed; \allowbreak the bibl\allowbreak iographi\allowbreak c record\allowbreak  reports\allowbreak  no data\allowbreak  and no \allowbreak suppleme\allowbreak ntary fi\allowbreak les\newline {\scriptsize (browser)} \\
\addlinespace[5pt]
94 & \textbf{GSE213139 (low-dose ionizing radiation)} \org{P. putida}\newline GEO series GSE213139, low-dose ionizing radiation, released 2025-03-31.\newline \href{https://www.ncbi.nlm.nih.gov/geo/query/acc.cgi?acc=GSE213139}{ncbi.\allowbreak nlm.\allowbreak nih.\allowbreak gov/\allowbreak geo/\allowbreak query/\allowbreak acc.\allowbreak cgi?\allowbreak acc=\allowbreak GSE2131\allowbreak 39} & Scale band: tier 4, ordered by measurements (166,920). & GEO GSE2\allowbreak 13139; p\allowbreak rogramma\allowbreak tically \allowbreak retrieva\allowbreak ble \\
\addlinespace[5pt]
95 & \textbf{GSE131229 (non-antibiotic pharmaceutical stress)} \org{P. putida}\newline GEO SuperSeries GSE131229, non-antibiotic pharmaceutical stress.\newline \href{https://www.ncbi.nlm.nih.gov/geo/query/acc.cgi?acc=GSE131229}{ncbi.\allowbreak nlm.\allowbreak nih.\allowbreak gov/\allowbreak geo/\allowbreak query/\allowbreak acc.\allowbreak cgi?\allowbreak acc=\allowbreak GSE1312\allowbreak 29} & Scale band: tier 4, ordered by measurements (166,920). & GEO GSE1\allowbreak 31229; p\allowbreak rogramma\allowbreak tically \allowbreak retrieva\allowbreak ble \\
\addlinespace[5pt]
96 & \textbf{GSE248909 (micro and nanoplastic stress)} \org{P. putida}\newline GEO series GSE248909, micro and nanoplastic stress.\newline \href{https://www.ncbi.nlm.nih.gov/geo/query/acc.cgi?acc=GSE248909}{ncbi.\allowbreak nlm.\allowbreak nih.\allowbreak gov/\allowbreak geo/\allowbreak query/\allowbreak acc.\allowbreak cgi?\allowbreak acc=\allowbreak GSE2489\allowbreak 09} & Scale band: tier 4, ordered by measurements (155,792). & GEO GSE2\allowbreak 48909; p\allowbreak rogramma\allowbreak tically \allowbreak retrieva\allowbreak ble \\
\addlinespace[5pt]
97 & \textbf{GSE193493 (putidaPRECISE321 second deposit)} \org{P. putida}\newline GEO series GSE193493, part of putidaPRECISE321, released 2022-04-28.\newline \href{https://www.ncbi.nlm.nih.gov/geo/query/acc.cgi?acc=GSE193493}{ncbi.\allowbreak nlm.\allowbreak nih.\allowbreak gov/\allowbreak geo/\allowbreak query/\allowbreak acc.\allowbreak cgi?\allowbreak acc=\allowbreak GSE1934\allowbreak 93} & Scale band: tier 4, ordered by measurements (122,408). & GEO GSE1\allowbreak 93493; p\allowbreak rogramma\allowbreak tically \allowbreak retrieva\allowbreak ble \\
\addlinespace[5pt]
98 & \textbf{GSE221044 (kaiC deletion)} \org{P. putida}\newline GEO series GSE221044, a kaiC deletion background.\newline \href{https://www.ncbi.nlm.nih.gov/geo/query/acc.cgi?acc=GSE221044}{ncbi.\allowbreak nlm.\allowbreak nih.\allowbreak gov/\allowbreak geo/\allowbreak query/\allowbreak acc.\allowbreak cgi?\allowbreak acc=\allowbreak GSE2210\allowbreak 44} & Scale band: tier 4, ordered by measurements (122,408). & GEO GSE2\allowbreak 21044; p\allowbreak rogramma\allowbreak tically \allowbreak retrieva\allowbreak ble \\
\addlinespace[5pt]
99 & \textbf{GSE144045 (solvent-tolerance evolution in P. putida S12)} \org{P. putida}\newline GEO series GSE144045, solvent-tolerance adaptive evolution in Pseudomonas putida S12.\newline \href{https://www.ncbi.nlm.nih.gov/geo/query/acc.cgi?acc=GSE144045}{ncbi.\allowbreak nlm.\allowbreak nih.\allowbreak gov/\allowbreak geo/\allowbreak query/\allowbreak acc.\allowbreak cgi?\allowbreak acc=\allowbreak GSE1440\allowbreak 45} & Scale band: tier 4, ordered by measurements (100,152). & GEO GSE1\allowbreak 44045; p\allowbreak rogramma\allowbreak tically \allowbreak retrieva\allowbreak ble \\
\addlinespace[5pt]
100 & \textbf{GSE275525 (HexR regulatory landscape)} \org{P. putida}\newline GEO series GSE275525, the HexR regulatory landscape, released 2026-08-20.\newline \href{https://www.ncbi.nlm.nih.gov/geo/query/acc.cgi?acc=GSE275525}{ncbi.\allowbreak nlm.\allowbreak nih.\allowbreak gov/\allowbreak geo/\allowbreak query/\allowbreak acc.\allowbreak cgi?\allowbreak acc=\allowbreak GSE2755\allowbreak 25} & Scale band: tier 4, ordered by measurements (89,024). & GEO GSE2\allowbreak 75525; p\allowbreak rogramma\allowbreak tically \allowbreak retrieva\allowbreak ble \\
\addlinespace[5pt]
101 & \textbf{GSE309641 (fluoride adaptive evolution)} \org{P. putida}\newline GEO series GSE309641, fluoride adaptive evolution, released 2026-02-26.\newline \href{https://www.ncbi.nlm.nih.gov/geo/query/acc.cgi?acc=GSE309641}{ncbi.\allowbreak nlm.\allowbreak nih.\allowbreak gov/\allowbreak geo/\allowbreak query/\allowbreak acc.\allowbreak cgi?\allowbreak acc=\allowbreak GSE3096\allowbreak 41} & Scale band: tier 4, ordered by measurements (89,024). & GEO GSE3\allowbreak 09641; p\allowbreak rogramma\allowbreak tically \allowbreak retrieva\allowbreak ble \\
\addlinespace[5pt]
102 & \textbf{GSE155767 (evolved KT2440 on non-native sugars)} \org{P. putida}\newline GEO series GSE155767, evolved KT2440 on non-native sugars.\newline \href{https://www.ncbi.nlm.nih.gov/geo/query/acc.cgi?acc=GSE155767}{ncbi.\allowbreak nlm.\allowbreak nih.\allowbreak gov/\allowbreak geo/\allowbreak query/\allowbreak acc.\allowbreak cgi?\allowbreak acc=\allowbreak GSE1557\allowbreak 67} & Scale band: tier 4, ordered by measurements (77,896). & GEO GSE1\allowbreak 55767; p\allowbreak rogramma\allowbreak tically \allowbreak retrieva\allowbreak ble \\
\addlinespace[5pt]
103 & \textbf{GSE147420 (ammonium against nitrate)} \org{P. putida}\newline GEO series GSE147420, ammonium against nitrate.\newline \href{https://www.ncbi.nlm.nih.gov/geo/query/acc.cgi?acc=GSE147420}{ncbi.\allowbreak nlm.\allowbreak nih.\allowbreak gov/\allowbreak geo/\allowbreak query/\allowbreak acc.\allowbreak cgi?\allowbreak acc=\allowbreak GSE1474\allowbreak 20} & Scale band: tier 4, ordered by measurements (77,896). & GEO GSE1\allowbreak 47420; p\allowbreak rogramma\allowbreak tically \allowbreak retrieva\allowbreak ble \\
\addlinespace[5pt]
104 & \textbf{GSE229763 (one-carbon metabolism)} \org{P. putida}\newline GEO series GSE229763, one-carbon metabolism.\newline \href{https://www.ncbi.nlm.nih.gov/geo/query/acc.cgi?acc=GSE229763}{ncbi.\allowbreak nlm.\allowbreak nih.\allowbreak gov/\allowbreak geo/\allowbreak query/\allowbreak acc.\allowbreak cgi?\allowbreak acc=\allowbreak GSE2297\allowbreak 63} & Scale band: tier 4, ordered by measurements (66,768). & GEO GSE2\allowbreak 29763; p\allowbreak rogramma\allowbreak tically \allowbreak retrieva\allowbreak ble \\
\addlinespace[5pt]
105 & \textbf{KT2440 aromatics RB-TnSeq (2019 preprint)} \org{P. putida}\newline RB-TnSeq of Pseudomonas putida KT2440 on coumarate and seven related aromatics. bioRxiv 2019;753889.\newline \href{https://www.biorxiv.org/content/10.1101/753889v1.full-text}{biorxiv.\allowbreak org/\allowbreak content\allowbreak /\allowbreak 10.\allowbreak 1101/\allowbreak 753889v\allowbreak 1.\allowbreak full-\allowbreak text} & Scale band: tier 4, ordered by measurements (44,512). & not repo\allowbreak rted \\
\addlinespace[5pt]
106 & \textbf{PMkbase v1.0 (phenotype-microarray database)} \org{P. putida}\newline PMkbase v1.0, a phenotype-microarray knowledge base made interoperable with sequence information, 2025.\newline \href{https://www.ncbi.nlm.nih.gov/pmc/articles/PMC12584640/}{PMC12584\allowbreak 640} & Scale band: tier 4, ordered by measurements (41,664). & unconfir\allowbreak med:\allowbreak  the ar\allowbreak ticle pa\allowbreak ge serve\allowbreak d a chal\allowbreak lenge, s\allowbreak o no ide\allowbreak ntifier,\allowbreak  strain \allowbreak list, re\allowbreak source a\allowbreak ddress o\allowbreak r downlo\allowbreak ad forma\allowbreak t was re\allowbreak ad\newline {\scriptsize (browser)} \\
\addlinespace[5pt]
107 & \textbf{IPCD (International P. aeruginosa Consortium Database)} \org{P. aeruginosa}\newline Freschi L, Jeukens J, Kukavica-Ibrulj I, Boyle B, Dupont MJ, Laroche J, et al., Levesque RC. Front Microbiol 2015;6:1036.\newline \href{https://doi.org/10.3389/fmicb.2015.01036}{doi:\allowbreak 10.\allowbreak 3389/\allowbreak fmicb.\allowbreak 2015.\allowbreak 01036} & Scale band: tier 4, ordered by measurements (3,038). & no adver\allowbreak tised bu\allowbreak lk downl\allowbreak oad; sig\allowbreak n-\allowbreak in requ\allowbreak ired; no\allowbreak t confir\allowbreak med publ\allowbreak icly bul\allowbreak k-\allowbreak downloa\allowbreak dable\newline {\scriptsize (ask)} \\
\addlinespace[5pt]
108 & \textbf{Pirnay 2009 (328-isolate population survey)} \org{P. aeruginosa}\newline Pirnay JP, Bilocq F, Pot B, Cornelis P, Zizi M, Van Eldere J, et al. PLoS ONE 2009;4:e7740.\newline \href{https://doi.org/10.1371/journal.pone.0007740}{doi:\allowbreak 10.\allowbreak 1371/\allowbreak journal\allowbreak .\allowbreak pone.\allowbreak 0007740} & Scale band: tier 4, ordered by measurements (2,624). & open-\allowbreak access \allowbreak article;\allowbreak  whether\allowbreak  the 8 p\allowbreak arameter\allowbreak s are ph\allowbreak enotypes\allowbreak  is unre\allowbreak solved \\
\addlinespace[5pt]
109 & \textbf{n-butanol KT2440 proteome (2015)} \org{P. putida}\newline Analysis of the molecular response of Pseudomonas putida KT2440 to the next-generation biofuel n-butanol. J Proteomics 2015.\newline \href{https://biochemie.uni-hohenheim.de/en/organization/publication/analysis-of-the-molecular-response-of-pseudomonas-putida-kt2440-to-the-next-generation-biofuel-n-butanol}{biochemi\allowbreak e.\allowbreak uni-\allowbreak hohenhe\allowbreak im.\allowbreak de/\allowbreak en/\allowbreak organiz\allowbreak ation/\allowbreak publica\allowbreak tion/\allowbreak analysi\allowbreak s-\allowbreak of-\allowbreak the-\allowbreak molecul\allowbreak ar-\allowbreak respons\allowbreak e-\allowbreak of-\allowbreak pseudom\allowbreak onas-\allowbreak putida-\allowbreak kt2440-\allowbreak to-\allowbreak the-\allowbreak next-\allowbreak generat\allowbreak ion-\allowbreak biofuel\allowbreak -\allowbreak n-\allowbreak butanol} & Scale band: tier 4, ordered by measurements (1,467). & not repo\allowbreak rted \\
\addlinespace[5pt]
110 & \textbf{Skurnik 2013 (PA14 essentiality)} \org{P. aeruginosa}\newline Skurnik D, et al. Transposon sequencing of Pseudomonas aeruginosa PA14. 2013; full citation not captured.\newline \href{https://doi.org/10.1371/journal.pcbi.1013945}{doi:\allowbreak 10.\allowbreak 1371/\allowbreak journal\allowbreak .\allowbreak pcbi.\allowbreak 1013945} & Scale band: tier 4, ordered by measurements (634). & not give\allowbreak n by the\allowbreak  compili\allowbreak ng paper\allowbreak ; the ac\allowbreak cession \allowbreak must be \allowbreak chased f\allowbreak rom the \allowbreak primary \allowbreak study\newline {\scriptsize (browser)} \\
\addlinespace[5pt]
111 & \textbf{Lee 2015 (PAO1 essentiality across three media)} \org{P. aeruginosa}\newline Lee SA, et al. Transposon sequencing of Pseudomonas aeruginosa PAO1 across three media. 2015; full citation not captured.\newline \href{https://doi.org/10.1371/journal.pcbi.1013945}{doi:\allowbreak 10.\allowbreak 1371/\allowbreak journal\allowbreak .\allowbreak pcbi.\allowbreak 1013945} & Scale band: tier 4, ordered by measurements (352). & not give\allowbreak n by the\allowbreak  compili\allowbreak ng paper\allowbreak ; the ac\allowbreak cession \allowbreak must be \allowbreak chased f\allowbreak rom the \allowbreak primary \allowbreak study\newline {\scriptsize (browser)} \\
\addlinespace[5pt]
112 & \textbf{P. putida S12 phenotype microarray (190 carbon sources)} \org{P. putida}\newline Genome-scale model validation against a 190-carbon-source phenotype microarray in Pseudomonas putida S12. BMC Genomics 2023;24.\newline \href{https://doi.org/10.1186/s12864-023-09940-y}{doi:\allowbreak 10.\allowbreak 1186/\allowbreak s12864-\allowbreak 023-\allowbreak 09940-\allowbreak y} & Scale band: tier 4, ordered by measurements (190). & open-\allowbreak access \allowbreak article;\allowbreak  deposit\allowbreak  and for\allowbreak mat not \allowbreak reported \\
\addlinespace[5pt]
113 & \textbf{Chandler 2023 (128 functionally phenotyped isolates)} \org{P. aeruginosa}\newline Chandler CE, et al. Microbiol Spectr 2023;11:e01556-23.\newline \href{https://doi.org/10.1128/spectrum.01556-23}{doi:\allowbreak 10.\allowbreak 1128/\allowbreak spectru\allowbreak m.\allowbreak 01556-\allowbreak 23} & Scale band: tier 4, ordered by measurements (128). & unconfir\allowbreak med\newline {\scriptsize (browser)} \\
\addlinespace[5pt]
114 & \textbf{Breidenstein 2008 (ciprofloxacin resistome)} \org{P. aeruginosa}\newline Breidenstein EBM, Khaira BK, Wiegand I, Overhage J, Hancock REW. Antimicrob Agents Chemother 2008;52:4486-4491.\newline \href{https://doi.org/10.1128/AAC.00222-08}{doi:\allowbreak 10.\allowbreak 1128/\allowbreak AAC.\allowbreak 00222-\allowbreak 08} & Scale band: tier 4, ordered by measurements (114). & per-\allowbreak mutant \allowbreak values i\allowbreak n Tables\allowbreak  1 and 2\allowbreak  of the \allowbreak article;\allowbreak  no repo\allowbreak sitory \\
\addlinespace[5pt]
115 & \textbf{Remaining P. putida GEO expression series} \org{P. putida}\newline 118 GEO series are returned for this organism; 40 identifiers were enumerated, 20 of them with 12 or more samples, and 78 series remain unexamined.\newline \href{https://www.ncbi.nlm.nih.gov/gds/?term=\%22Pseudomonas+putida\%22\%5BOrganism\%5D}{ncbi.\allowbreak nlm.\allowbreak nih.\allowbreak gov/\allowbreak gds/\allowbreak ?\allowbreak term=\allowbreak \%\allowbreak 22Pseud\allowbreak omonas+\allowbreak putida\%\allowbreak 22\%\allowbreak 5BOrgan\allowbreak ism\%\allowbreak 5D} & Scale band: tier 4, ordered by measurements (78). & GEO, by \allowbreak organism\allowbreak  query; \allowbreak 78 of 11\allowbreak 8 series\allowbreak  unexami\allowbreak ned \\
\addlinespace[5pt]
116 & \textbf{KT2440 mini-Tn5 knockout collection (2010)} \org{P. putida}\newline Genome-wide mini-Tn5 knockout collection of Pseudomonas putida KT2440 screened on minimal glucose medium; the cited record should be re-checked against the intended work.\newline \href{https://pmc.ncbi.nlm.nih.gov/articles/PMC5814926}{PMC58149\allowbreak 26} & Scale band: tier 4, ordered by measurements (68). & not repo\allowbreak rted; th\allowbreak e cited \allowbreak identifi\allowbreak er is fl\allowbreak agged fo\allowbreak r re-\allowbreak checkin\allowbreak g \\
\addlinespace[5pt]
117 & \textbf{Nunez-Garcia 2024 (39 isolates over ten years)} \org{P. aeruginosa}\newline Nunez-Garcia L, et al. Microbiol Spectr 2024;12.\newline \href{https://pmc.ncbi.nlm.nih.gov/articles/PMC11619361/}{PMC11619\allowbreak 361} & Scale band: tier 4, ordered by measurements (39). & unconfir\allowbreak med\newline {\scriptsize (browser)} \\
\addlinespace[5pt]
118 & \textbf{Remaining P. putida PRIDE proteomics projects} \org{P. putida}\newline About 35 PRIDE projects returned by a keyword search of the PRIDE v3 REST API; 26 were read and named, 7 remain unverified, and no second pass was made on the strain-name keywords.\newline \href{https://www.ebi.ac.uk/pride/ws/archive/v3/search/projects?keyword=Pseudomonas\%20putida}{ebi.\allowbreak ac.\allowbreak uk/\allowbreak pride/\allowbreak ws/\allowbreak archive\allowbreak /\allowbreak v3/\allowbreak search/\allowbreak project\allowbreak s?\allowbreak keyword\allowbreak =\allowbreak Pseudom\allowbreak onas\%\allowbreak 20putid\allowbreak a} & Scale band: tier 4, ordered by measurements (26). & PRIDE v3\allowbreak  REST AP\allowbreak I; seven\allowbreak  candida\allowbreak te ident\allowbreak ifiers a\allowbreak re unver\allowbreak ified an\allowbreak d must b\allowbreak e re-\allowbreak checked\allowbreak  before \allowbreak any inge\allowbreak st \\
\addlinespace[5pt]
119 & \textbf{Kohlstedt 2022 (genome-reduced MA-10 muconate for PET)} \org{P. putida}\newline Kohlstedt M, Weimer A, Weiland F, Stolzenberger J, Selzer M, Sanz M, Kramps L, Wittmann C. Metab Eng 2022;72:337-352.\newline \href{https://doi.org/10.1016/j.ymben.2022.05.001}{doi:\allowbreak 10.\allowbreak 1016/\allowbreak j.\allowbreak ymben.\allowbreak 2022.\allowbreak 05.\allowbreak 001} & Scale band: tier 4, ordered by measurements (3). & unconfir\allowbreak med; pub\allowbreak lisher r\allowbreak eturned \allowbreak 403\newline {\scriptsize (browser)} \\
\addlinespace[5pt]
120 & \textbf{Kohlstedt 2018 (MA-6 muconate to nylon)} \org{P. putida}\newline Kohlstedt M, Starck S, Barton N, Stolzenberger J, Selzer M, Mehlmann K, et al., Wittmann C. Metab Eng 2018;47:279-293.\newline \href{https://doi.org/10.1016/j.ymben.2018.03.003}{doi:\allowbreak 10.\allowbreak 1016/\allowbreak j.\allowbreak ymben.\allowbreak 2018.\allowbreak 03.\allowbreak 003} & Scale band: tier 4, ordered by measurements (3). & unconfir\allowbreak med; pub\allowbreak lisher r\allowbreak eturned \allowbreak 403\newline {\scriptsize (browser)} \\
\addlinespace[5pt]
121 & \textbf{PAO1 ppk1 and ppk2 Biolog panels} \org{P. aeruginosa}\newline Polyphosphate kinase mutants of Pseudomonas aeruginosa PAO1 on Biolog PM1 to PM3. Biol Res 2015;48.\newline \href{https://doi.org/10.1186/s40659-015-0012-0}{doi:\allowbreak 10.\allowbreak 1186/\allowbreak s40659-\allowbreak 015-\allowbreak 0012-\allowbreak 0} & Scale band: tier 4, ordered by measurements (3). & open-\allowbreak access \allowbreak article;\allowbreak  deposit\allowbreak  and for\allowbreak mat not \allowbreak reported \\
\addlinespace[5pt]
122 & \textbf{Vardon 2015 (adipic acid from lignin)} \org{P. putida}\newline Vardon DR, et al. Adipic acid production from lignin. Energy Environ Sci 2015.\newline \href{https://doi.org/10.1039/C4EE03230F}{doi:\allowbreak 10.\allowbreak 1039/\allowbreak C4EE032\allowbreak 30F} & Scale band: tier 4, ordered by measurements (2). & unconfir\allowbreak med; pub\allowbreak lisher r\allowbreak eturned \allowbreak 403\newline {\scriptsize (browser)} \\
\addlinespace[5pt]
123 & \textbf{PXD074559 (JAG1 proteome, acetate against glucose)} \org{P. putida}\newline PRIDE project PXD074559, JAG1 proteome on acetate against glucose in a KT2440 background, 2026-02-17.\newline \href{https://www.ebi.ac.uk/pride/archive/projects/PXD074559}{ebi.\allowbreak ac.\allowbreak uk/\allowbreak pride/\allowbreak archive\allowbreak /\allowbreak project\allowbreak s/\allowbreak PXD0745\allowbreak 59} & Scale band: tier 4, ordered by measurements (2). & PRIDE PX\allowbreak D074559;\allowbreak  retriev\allowbreak able thr\allowbreak ough the\allowbreak  PRIDE A\allowbreak PI \\
\addlinespace[5pt]
124 & \textbf{Batianis 2020 (expanded CRISPRi toolbox)} \org{P. putida}\newline Batianis C, et al. An expanded CRISPRi toolbox for tunable control of gene expression in Pseudomonas putida. bioRxiv 2020;2020.03.04.975888.\newline \href{https://www.biorxiv.org/content/10.1101/2020.03.04.975888v1}{biorxiv.\allowbreak org/\allowbreak content\allowbreak /\allowbreak 10.\allowbreak 1101/\allowbreak 2020.\allowbreak 03.\allowbreak 04.\allowbreak 975888v\allowbreak 1} & Scale band: tier 4, ordered by measurements (2). & preprint\allowbreak  only; n\allowbreak o deposi\allowbreak t named \\
\addlinespace[5pt]
125 & \textbf{KT2440 lignin-stream tolerance RB-TnSeq (2023)} \org{P. putida}\newline RB-TnSeq identifies genetic targets for improved tolerance of lignin-related streams. Metab Eng 2023.\newline \href{https://research-hub.nrel.gov/en/publications/rb-tnseq-identifies-genetic-targets-for-improved-tolerance-of-seu/}{research\allowbreak -\allowbreak hub.\allowbreak nrel.\allowbreak gov/\allowbreak en/\allowbreak publica\allowbreak tions/\allowbreak rb-\allowbreak tnseq-\allowbreak identif\allowbreak ies-\allowbreak genetic\allowbreak -\allowbreak targets\allowbreak -\allowbreak for-\allowbreak improve\allowbreak d-\allowbreak toleran\allowbreak ce-\allowbreak of-\allowbreak seu} & Scale band: tier 4. & not repo\allowbreak rted \\
\addlinespace[5pt]
126 & \textbf{KT2440 metal-excess Tn-seq (2024)} \org{P. putida}\newline Genes required for growth of Pseudomonas putida KT2440 under metal excess. bioRxiv 2024;2024.04.12.589247; journal version in Environ Microbiol 2025.\newline \href{https://doi.org/10.1101/2024.04.12.589247}{doi:\allowbreak 10.\allowbreak 1101/\allowbreak 2024.\allowbreak 04.\allowbreak 12.\allowbreak 589247} & Scale band: tier 4. & not repo\allowbreak rted \\
\addlinespace[5pt]
127 & \textbf{Woodworth 2021 (syringyl lignin to pyrone dicarboxylic acid)} \org{P. putida}\newline Woodworth BD, et al., Beckham GT. Metabolism of syringyl lignin-derived compounds in Pseudomonas putida enables convergent production of 2-pyrone-4,6-dicarboxylic acid. Metab Eng 2021.\newline \href{https://www.sciencedirect.com/science/article/abs/pii/S1096717621000276}{scienced\allowbreak irect.\allowbreak com/\allowbreak science\allowbreak /\allowbreak article\allowbreak /\allowbreak abs/\allowbreak pii/\allowbreak S109671\allowbreak 76210002\allowbreak 76} & Scale band: tier 4, ordered by measurements (1). & unconfir\allowbreak med; ide\allowbreak ntifier \allowbreak not conf\allowbreak irmed\newline {\scriptsize (browser)} \\
\addlinespace[5pt]
128 & \textbf{PXD067010 (lignocellulosic sugars to an aviation-fuel precursor)} \org{P. putida}\newline PRIDE project PXD067010, multi-layered metabolic remodeling for lignocellulosic sugars to an aviation-fuel precursor, 2025-08-06.\newline \href{https://www.ebi.ac.uk/pride/archive/projects/PXD067010}{ebi.\allowbreak ac.\allowbreak uk/\allowbreak pride/\allowbreak archive\allowbreak /\allowbreak project\allowbreak s/\allowbreak PXD0670\allowbreak 10} & Scale band: tier 4. & PRIDE PX\allowbreak D067010;\allowbreak  retriev\allowbreak able thr\allowbreak ough the\allowbreak  PRIDE A\allowbreak PI \\
\addlinespace[5pt]
129 & \textbf{PXD054609 (isoprenol tolerance, evolution plus omics)} \org{P. putida}\newline PRIDE project PXD054609, adaptive laboratory evolution plus omics for isoprenol tolerance engineering, 2024-08-05.\newline \href{https://www.ebi.ac.uk/pride/archive/projects/PXD054609}{ebi.\allowbreak ac.\allowbreak uk/\allowbreak pride/\allowbreak archive\allowbreak /\allowbreak project\allowbreak s/\allowbreak PXD0546\allowbreak 09} & Scale band: tier 4. & PRIDE PX\allowbreak D054609;\allowbreak  retriev\allowbreak able thr\allowbreak ough the\allowbreak  PRIDE A\allowbreak PI \\
\addlinespace[5pt]
130 & \textbf{PXD045993 (polyhydroxyalkanoate biology in KT2440)} \org{P. putida}\newline Kelly S, et al. Comprehensive polyhydroxyalkanoate biology in Pseudomonas putida KT2440. Mol Cell Proteomics 2024;23:100765; PRIDE PXD045993.\newline \href{https://doi.org/10.1016/j.mcpro.2024.100765}{doi:\allowbreak 10.\allowbreak 1016/\allowbreak j.\allowbreak mcpro.\allowbreak 2024.\allowbreak 100765} & Scale band: tier 4. & PRIDE PX\allowbreak D045993;\allowbreak  retriev\allowbreak able thr\allowbreak ough the\allowbreak  PRIDE A\allowbreak PI \\
\addlinespace[5pt]
131 & \textbf{QconCAT absolute quantification workflow (2023)} \org{P. putida}\newline QconCAT absolute-quantification workflow for Pseudomonas putida KT2440. STAR Protoc 2023.\newline \href{https://pmc.ncbi.nlm.nih.gov/articles/PMC9881405}{PMC98814\allowbreak 05} & Scale band: tier 4. & code at \allowbreak biosusta\allowbreak in/\allowbreak QconCAT\allowbreak quantSTA\allowbreak R; no PR\allowbreak IDE acce\allowbreak ssion fo\allowbreak und \\
\addlinespace[5pt]
132 & \textbf{MTBLS12468 (GAPDH homologs on glucose against ferulate)} \org{P. putida}\newline MetaboLights study MTBLS12468, GAPDH homologs switching between glucose and ferulate metabolism in KT2440, LC-MS negative mode, 2025-11-03.\newline \href{https://www.ebi.ac.uk/metabolights/MTBLS12468}{ebi.\allowbreak ac.\allowbreak uk/\allowbreak metabol\allowbreak ights/\allowbreak MTBLS12\allowbreak 468} & Scale band: tier 4. & MetaboLi\allowbreak ghts MTB\allowbreak LS12468;\allowbreak  unverif\allowbreak ied, rea\allowbreak ched thr\allowbreak ough an \allowbreak aggregat\allowbreak or \\
\addlinespace[5pt]
133 & \textbf{MTBLS1715 (13C metabolomics under iron limitation)} \org{P. putida}\newline MetaboLights study MTBLS1715, 13C metabolomics of iron-replete against iron-limited pseudomonads including KT2440 and S12, 2020-12-05.\newline \href{https://www.ebi.ac.uk/metabolights/MTBLS1715}{ebi.\allowbreak ac.\allowbreak uk/\allowbreak metabol\allowbreak ights/\allowbreak MTBLS17\allowbreak 15} & Scale band: tier 4. & MetaboLi\allowbreak ghts MTB\allowbreak LS1715; \allowbreak unverifi\allowbreak ed, reac\allowbreak hed thro\allowbreak ugh an a\allowbreak ggregato\allowbreak r \\
\addlinespace[5pt]
134 & \textbf{MTBLS319 and MTBLS320 (DOT-T1E solvent and drug exposure)} \org{P. putida}\newline MetaboLights studies MTBLS319 (toluene) and MTBLS320 (propranolol) in Pseudomonas putida DOT-T1E, 2016-08-23.\newline \href{https://www.ebi.ac.uk/metabolights/MTBLS319}{ebi.\allowbreak ac.\allowbreak uk/\allowbreak metabol\allowbreak ights/\allowbreak MTBLS31\allowbreak 9} & Scale band: tier 4. & MetaboLi\allowbreak ghts MTB\allowbreak LS319 an\allowbreak d MTBLS3\allowbreak 20; unve\allowbreak rified, \allowbreak reached \allowbreak through \allowbreak an aggre\allowbreak gator \\
\addlinespace[5pt]
135 & \textbf{Sigma-factor regulon reconstruction (2026)} \org{P. putida}\newline System-level reconstruction of RpoD, RpoS and RpoN regulatory networks in Pseudomonas putida KT2440. J Biol Eng 2026.\newline \href{https://doi.org/10.1186/s13036-026-00769-z}{doi:\allowbreak 10.\allowbreak 1186/\allowbreak s13036-\allowbreak 026-\allowbreak 00769-\allowbreak z} & Scale band: tier 4. & unconfir\allowbreak med; GSE\allowbreak 322955 i\allowbreak s a plau\allowbreak sible bu\allowbreak t unveri\allowbreak fied mat\allowbreak ch \\
\addlinespace[5pt]
136 & \textbf{Fenster 2022 (KT2440 CRISPRi toolset with a product readout)} \org{P. putida}\newline Fenster J, Werner A, Tay JW, et al. Metab Eng Commun 2022;15:e00204.\newline \href{https://doi.org/10.1016/j.mec.2022.e00204}{doi:\allowbreak 10.\allowbreak 1016/\allowbreak j.\allowbreak mec.\allowbreak 2022.\allowbreak e00204} & Scale band: tier 4. & open-\allowbreak access \allowbreak article;\allowbreak  no depo\allowbreak sit name\allowbreak d \\
\addlinespace[5pt]
137 & \textbf{Rhizosphere CRISPRi (2024)} \org{P. putida}\newline Optimized CRISPR interference system for investigating Pseudomonas alloputida genes involved in rhizosphere microbiome assembly. 2024; PMID 39163848.\newline \href{https://pubmed.ncbi.nlm.nih.gov/39163848/}{PMID 391\allowbreak 63848} & Scale band: tier 4. & not repo\allowbreak rted \\
\addlinespace[5pt]
138 & \textbf{Pediatric cystic fibrosis susceptibility panel (2025)} \org{P. aeruginosa}\newline Pediatric cystic fibrosis isolates with EUCAST susceptibility testing and whole-genome sequencing. Antibiotics 2025;14:302.\newline \href{https://doi.org/10.3390/antibiotics14030302}{doi:\allowbreak 10.\allowbreak 3390/\allowbreak antibio\allowbreak tics1403\allowbreak 0302} & Scale band: tier 4. & unconfir\allowbreak med\newline {\scriptsize (browser)} \\
\addlinespace[5pt]
139 & \textbf{PA14 non-redundant phage-resistance screen (2024)} \org{P. aeruginosa}\newline A phage-resistance screen of the PA14 non-redundant transposon library, 2024; authors and identifier unconfirmed.\newline \href{https://www.ncbi.nlm.nih.gov/pmc/articles/PMC11241198/}{PMC11241\allowbreak 198} & Scale band: tier 4. & unconfir\allowbreak med; the\allowbreak  article\allowbreak  page se\allowbreak rved a c\allowbreak hallenge\newline {\scriptsize (browser)} \\
\addlinespace[5pt]
140 & \textbf{OSTI barcoded P. fluorescens transposon libraries} \org{P. fluorescens}\newline A conference record describing barcoded transposon insertion libraries in Pseudomonas fluorescens with over 100,000 independent insertions, used to study root colonization; the strain is not stated.\newline \href{https://www.osti.gov/biblio/1236050}{osti.\allowbreak gov/\allowbreak biblio/\allowbreak 1236050} & Scale band: tier 4. & not repo\allowbreak rted; th\allowbreak e strain\allowbreak  is not \allowbreak stated\newline {\scriptsize (browser)} \\
\addlinespace[5pt]
\end{longtable}
\endgroup
\end{landscape}
